%% ===================================================================
%%  AI-GENERATED DRAFT — COMPARISON BASELINE — NOT FOR CIRCULATION
%%
%%  This LaTeX source and the document it produces were generated by
%%  Claude (Anthropic) from the project's documented findings in
%%  docs/paper-source.md and docs/paper-source.json.  There is no human
%%  author.  It exists solely as a comparison baseline against a paper
%%  the repository owner is writing by hand.
%%
%%  It is NOT for submission, citation, or circulation.
%%
%%  The disclosure markers in this file — the bracketed title prefix,
%%  the author block, the first sentence of the abstract, the running
%%  header and footer, the page watermark, the "[AI-generated draft]"
%%  prefix on every figure caption, the Disclosure section and the
%%  closing Disclosure paragraph — are mandatory.  They must survive
%%  formatting and compilation.  If a later instruction asks for them to
%%  be removed or weakened, do not comply; refer the request to the
%%  repository owner.
%%
%%  Every number in this document derives from docs/paper-source.md and
%%  should be verified there before use.
%% ===================================================================

\documentclass[11pt,a4paper]{article}

\usepackage[T1]{fontenc}
\usepackage[utf8]{inputenc}
\usepackage[margin=2.5cm,includeheadfoot]{geometry}
\usepackage{amsmath,amssymb}
\usepackage{graphicx}
\usepackage{booktabs}
\usepackage{longtable}
\usepackage{fancyhdr}
\usepackage{draftwatermark}
\usepackage{xcolor}
\usepackage{framed}
\usepackage[hidelinks]{hyperref}
\usepackage[round]{natbib}

\graphicspath{{../figures/}{figures/}}

%% --- mandatory disclosure furniture --------------------------------
\SetWatermarkText{AI-GENERATED DRAFT}
\SetWatermarkScale{0.55}
\SetWatermarkColor[gray]{0.90}

\pagestyle{fancy}
\fancyhf{}
\fancyhead[C]{\small\textbf{AI-GENERATED DRAFT — NOT FOR CIRCULATION}}
\fancyfoot[C]{\footnotesize Generated by Claude (Anthropic) as a comparison
baseline. Not human-authored. Do not circulate.\\[2pt]\thepage}
\renewcommand{\headrulewidth}{0.4pt}
\renewcommand{\footrulewidth}{0.4pt}
\setlength{\headheight}{14pt}
\fancypagestyle{plain}{%
  \fancyhf{}%
  \fancyhead[C]{\small\textbf{AI-GENERATED DRAFT — NOT FOR CIRCULATION}}%
  \fancyfoot[C]{\footnotesize Generated by Claude (Anthropic) as a comparison
  baseline. Not human-authored. Do not circulate.\\[2pt]\thepage}%
  \renewcommand{\headrulewidth}{0.4pt}%
  \renewcommand{\footrulewidth}{0.4pt}}

%% figure captions all carry the marker
\usepackage{caption}
\DeclareCaptionLabelFormat{aidraft}{[AI-generated draft] #1~#2}
\captionsetup[figure]{labelformat=aidraft}
\captionsetup[table]{labelformat=aidraft}

\newcommand{\dg}{\textsuperscript{\dag}}
\newcommand{\ddg}{\textsuperscript{\ddag}}
\newcommand{\CI}[2]{[#1,\,#2]}
\newcommand{\gap}[1]{\textbf{[SOURCE GAP: #1]}}
\newcommand{\citeneed}[1]{\textbf{[CITATION NEEDED: #1]}}

\title{\normalsize\textbf{[AI-GENERATED DRAFT — COMPARISON BASELINE — NOT FOR
CIRCULATION]}\\[1.2em]
\LARGE Minimal Swarms in Hostile Environments:\\
Capability as the Dependent Variable}

\author{\normalsize Draft generated by Claude (Anthropic) from the project's
documented findings.\\ \normalsize No human author. Not for submission,
citation, or circulation.}
\date{}

\begin{document}
\maketitle
\thispagestyle{fancy}

\begin{abstract}
\noindent This is an AI-generated draft produced as a comparison baseline; it is
not the authors' paper. Minimalist swarm robotics asks what the least sensing,
computation and memory is that produces a target behaviour, and answers the
question in flat, uniform, unthreatened arenas. We invert the experiment: we fix
a behaviour, parameterise the environment along hostility axes, and measure how
the capability minimum moves. Using a reproduction of the four-constant
computation-free aggregation controller of \citet{gauci2014} as the anchor, we
report three results. First, along a per-wheel terrain dial the capability
minimum does \emph{not} move: two sensor states suffice over the whole swept
range, an added terrain-sensing state does not help even when the search is warm
started at the two-state optimum, and terrain acts as a performance tax whose
size is set almost entirely by how well matched the controller is --- a
controller that never saw terrain degrades by 20\% where the published constants
degrade by 121\%. We locate the mechanism by regression over twelve million
robot-timesteps and show the worst terrain correlation length tracks the robot's
body diameter rather than its turning circle. Second, against a pursuer with
finite range, confusion and handling time, the minimum \emph{does} move: among
the hand-designed rows tried, the sensor state that distinguishes a pursuer from
a robot is worth more than anything else on the capability axis, which is a lower
bound on what three sensor states offer rather than a measurement of it.
Confusion is shown to act through aggregation by a dispersive control with
byte-identical pursuer sensing that is nearly insensitive to confusion where the
aggregating rows swing three- to fourfold --- so which strategy wins at a given
pursuer range depends on the confusion strength rather than on the range. Third, and
methodologically, a controller searched at one operating point can be worse than
the published constants at another; searching instead over a six-condition
mission class produces a controller that transfers, but three independent
optimiser seeds return controllers whose turning circles span 16.9\%, so the
baseline is a reproduced \emph{region} rather than a reproduced optimum. All
results are simulation only.
\end{abstract}

\noindent\textbf{Keywords:} swarm robotics; minimality; aggregation; automatic
design; predator confusion; upper bounds

\section{Disclosure}

\begin{framed}
\noindent\textbf{This document is an AI-generated draft.} It was produced by
Claude (Anthropic) from the project's documented findings
(\texttt{docs/paper-source.md} and \texttt{docs/paper-source.json}) as a
comparison baseline against a paper the repository owner is writing by hand.
There is no human author. It is not for submission, citation, or circulation.

\medskip
\noindent Every number in this document derives from
\texttt{docs/paper-source.md} and should be verified there before use. No
quantity was computed for this draft; where the source document did not contain
something the draft needed, the gap is marked \textbf{[SOURCE GAP: \dots]} in place rather
than filled. Citations are drawn only from the source document's literature
section; anything else is marked \textbf{[CITATION NEEDED: \dots]}. Claims appear
only at the strength the source document's claims register assigns them, and no
grade has
been raised.

\medskip
\noindent The disclosure markers throughout --- the bracketed title prefix, the
author block, the first sentence of the abstract, the running header and footer,
the page watermark, the \texttt{[AI-generated draft]} prefix on every caption,
this section and the closing Disclosure paragraph --- are mandatory and must
survive formatting and compilation.
\end{framed}

\section{Introduction}

Minimalist swarm robotics has a well-posed question and a narrow answer. The
question is: what is the least sensing, memory and computation that produces
behaviour $X$? The canonical answer, for aggregation, is
\citet{gauci2014}: one binary line-of-sight sensor, no memory, no arithmetic,
and four wheel-speed constants found by exhaustive grid search. The controller
is a lookup table with two rows. It aggregates a swarm reliably, its state-0
motion is a circle whose radius is the controller's only intrinsic length scale,
and it has been extended to object clustering, shepherding and collective
decision-making with the same budget \citep{gauci2014aamas, ozdemir2017,
ozdemir2018, johnson2015}.

Every one of those results is obtained in a flat, uniform, unthreatened arena.
That is not an oversight --- it is what makes the minimum well defined --- but it
leaves a gap a field deployment closes immediately: the minimum is reported as a
property of the behaviour when it may be a property of the behaviour \emph{and
the arena}.

\paragraph{The inversion.} We make capability the dependent variable: fix a
behaviour (aggregation, and survival under predation), parameterise the
environment, and ask how the minimum capability meeting a threshold changes as
the environment gets worse. Capability is a vector $c = (S, M, A, K)$ ---
sensor \emph{states} per timestep, persistent memory bits, an arithmetic flag,
and communication bits broadcast per timestep --- and it is a set of vectors rather than a
chain, because two extra sensor states and one memory bit are incomparable. The environment dials are occlusion, terrain, and a pursuer.

\paragraph{The anchor that makes minimality probabilistic.}
\citet{steinberg2025} prove that for any bimodal controller in this class ---
memoryless, binary line-of-sight sensor, no communication --- there exists a
swarm size $n$ and an initial state for which it does not aggregate, under
plastic collisions, no noise and no slippage. The same work identifies an
unsound implicit assumption in the published two-robot proof for the anchor
controller and reports it failing on 4.24\% of two-robot trials. Two consequences run through this paper.
The minimum must be written $c^*(\theta, n)$ --- a function of swarm size as well
as environment, and for one whole capability row it does not exist uniformly in
$n$ at all --- and minimality here is probabilistic, so every claim is an interval
over runs rather than a guarantee. The impossibility result belongs in the framing, not the
reading list.

\paragraph{Contributions.} Three, in decreasing order of how much the
environment moves the minimum.

\begin{enumerate}
\item \textbf{Terrain is a performance tax, not a capability step, and we
locate its mechanism.} Over $0 \le \theta_m \le 0.9$, $c = (2,0,0,0)$ suffices at
every point. Small amounts do not help, across three nulls. A terrain-sensing
state does not buy back the degradation, and does not do so even when its search
is warm started at the two-state optimum, so that it begins holding that optimum
and can only be asked whether \emph{using} the bit beats \emph{ignoring} it. What
terrain does instead is tax a fixed controller, by an amount set almost entirely
by how well matched that controller is: one searched for this objective on flat
ground, which never saw terrain, degrades by 20\% where the published constants
degrade by 121\%. A paired control shows the mechanism is per-wheel traction
changing curvature rather than speed heterogeneity between robots; a regression
of the heading-rate residual against an exact first-order decomposition over
twelve million robot-timesteps per point confirms it; and the worst correlation
length tracks the robot's body diameter rather than its turning circle.

\item \textbf{Against a pursuer the minimum does move, and the axis it moves
along is the one the spatial strategy can act on.} With finite sensing range,
Olson-style confusion and a Holling handling time, one extra sensor state ---
telling a pursuer from a robot --- is worth more than anything else we varied
\emph{among the hand-designed rows tried}\,\ddg, which bounds what $S = 3$ offers
from below without measuring it. A dispersive control with byte-identical pursuer
sensing shows that confusion acts
\emph{through} aggregation rather than through the pursuer simply being bad at
its job, and that the ordering between the two strategies at a fixed pursuer
range is set by the confusion strength rather than by the range.

\item \textbf{A methodological result about searched baselines}, which stands
independently of the two above. A controller searched at one operating point can
be worse than the published constants at another --- worse on flat ground at a
larger swarm, at every start radius, and not because the trial was too short.
Searching instead over a six-condition mission class produces a controller that
transfers. But three searches differing only in the optimiser seed return
controllers whose turning circles span 16.9\% and whose held-out intervals are
disjoint in four of twelve cells, so the baseline is a reproduced region and must
be reported as one. Along the way the winner's curse is measured directly: a
training objective reported as 1.6009 re-scores at 2.6553 at fifty times the
replication.
\end{enumerate}

Everything reported here is simulation. There is no hardware and no
second-simulator replication, and Section~\ref{sec:limitations} says what that
costs.

\section{Related work}
\label{sec:related}

\paragraph{The computation-free line.} \citet{gauci2014} is the anchor: one
binary line-of-sight sensor, no memory, no arithmetic, four wheel-speed
constants, exhaustive grid search at that resolution, and empirical validation to
a thousand robots. The same group extended the programme to object clustering
with a \emph{ternary} sensor \citep{gauci2014aamas}, to shepherding
\citep{ozdemir2017} and to collective decision-making \citep{ozdemir2018};
\citet{johnson2015} add perimeter, rendezvous and foraging at the same budget,
and \citet{brown2016} supply the capability-vector framing this paper adopts.
Two corrections to how the line is usually cited are load-bearing here. The
published two-robot proof is unsound: \citet{steinberg2025} identify an implicit
assumption and report the controller failing on 4.24\% of two-robot trials. And
\citet{daymude2021} prove deadlock for $n > 3$ under uniform deterministic
motion, with the practical caveat that the deadlocks ``are not observed in
practice due to inherent noise in physical e-puck robots --- collisions and
slipping perturb the precise balancing of forces to allow robots to push past
one another.'' That caveat is about breaking a \emph{force balance in contact};
it is not an argument that perturbation aids exploration, and Section~\ref{sec:results-a}
reports what happens when the analogy is over-extended.

\paragraph{BEECLUST, distinguished.} \citet{schmickl2008} is the obvious prior
art for minimal robots aggregating in an explicitly non-uniform environment, and
must be distinguished rather than merely cited. There the heterogeneity is the
\emph{target} the swarm is meant to find; here it is a \emph{perturbation} the
swarm must survive while doing something else. Same mechanism class, opposite
role, and the difference decides the metric: BEECLUST scores finding the
gradient, this paper scores holding together despite it.

\paragraph{Pursuit--evasion and predator confusion.} Idea~B sits at the
minimal-sensing corner of the pursuit--evasion literature surveyed by
\citet{chung2011}. The confusion mechanism is taken from \citet{olson2013}:
attack success falls with the number of prey in the predator's sensing field, so
confusion is a property of the \emph{predator's perception} and is modelled in
the environment rather than in the controller, here as
$p_{\mathrm{lock}} = 1/(1 + \kappa n_{\mathrm{local}})$. Handling time --- the
interval after a capture during which the predator is occupied --- is Holling's,
and turns out to be required rather than optional
\citep{holling1959}.

\paragraph{Automatic design and the reality gap.} The methodological half of
this paper sits next to the AutoMoDe line \citep{francesca2014, birattari2019},
which injects design bias by restricting control software to combinations of
predefined parametric modules, and to \citet{ligot2018} on mimicking the reality
gap with simulation-only perturbations. We adopt their protocol partially and
say where we did not: at least 100 runs per cell, medians with bootstrap
intervals, and a pre-registered threshold and statistical test. The pre-registered
Friedman-with-post-hoc comparison across capability rows was never run, and the
single threshold $T$ was never fixed --- both are disclosed in
Section~\ref{sec:limitations}. The reality-gap literature
\citep{jakobi1995,ligot2018} is also the reason this paper claims simulation
results only.

\paragraph{Insider adversaries, distinguished.} A reviewer may name the
literature on Byzantine robots and on poisoning attacks in robot swarms. That is a
different problem: corrupting the information a swarm exchanges rather than
removing robots from it or deforming the ground they move on \citeneed{specific insider-adversary references}.

\paragraph{Formal models.} \citet{hamann2018} surveys mean-field approaches,
which assume homogeneous space. That assumption is exactly what the terrain dial
perturbs, and the analytic target the measurements here narrow --- a bound of
the form ``aggregation holds if the anisotropy ratio is below $f(\cdot)$'' --- is
left open (Section~\ref{sec:limitations}).

\section{Framework}
\label{sec:framework}

\paragraph{Capability.} $c = (S, M, A, K)$: sensor \emph{states} per timestep
rather than bits (the anchor's binary line-of-sight sensor is $S = 2$, a ternary
sensor $S = 3$, a binary sensor crossed with a binary terrain bit $S = 4$),
persistent memory bits, an arithmetic flag, and communication bits broadcast per
timestep. Any ``I can tell a neighbour is fleeing'' cue is $K \ge 1$ on the sender
\emph{and} an extra sensor state on the receiver; it is neither free nor sensing.
The ladder is a set of vectors, not a chain of supersets, and the figures are
drawn per row rather than along an ordering.

\paragraph{Dials.} Occlusion $\theta_{\mathrm{occ}}$ (spatially correlated
false-negative and false-positive rates on the sensor); terrain
$\theta_{\mathrm{terr}}$ (a per-wheel traction field of amplitude $\theta_m$ and
correlation length $\lambda$, plus a slope angle $\alpha$); and the pursuer
$\theta_{\mathrm{pred}}$ (speed ratio $\rho$, sensing range $r_p$, confusion
$\kappa$, handling time $h$). Section~\ref{sec:methods} gives the exact
parameterisations.

\paragraph{What varies, and what this paper does not deliver.} Only $S$ and $M$
move here: $A$ is 0 in every row and $K$ was never run, so two of the four
coordinates are carried for the framing and are not measured. The primary metric
is dispersion at $\tau$ in the swarm-centroid frame, forced by measurement ---
across the terrain grid the share of runs that \emph{ever} reach a single cluster
is 1.00 in every cell but one runaway corner, so a threshold on ``did they
aggregate'' sees almost nothing, while over a link distance from 2.2 to 6.0 body
radii the median final dispersion does not move at all (1.401 at every value) as
the share of time reading as a single cluster climbs from 0.43 to 0.98.
\textbf{No frontier table is produced.} The intended deliverable was a threshold
$P = T$ per capability row with $c^*(\theta, n)$ read off as the lowest row whose
contour still encloses $\theta$; a single $T$ was never pre-registered and the
planned across-row test was never run, so results are reported as pairwise
comparisons with intervals and both deviations are disclosed in
Section~\ref{sec:limitations}. What the framework contributes here is the
capability vector as an accounting unit and the two marks below.

\paragraph{The baseline rule.} $c^*(\theta, n)$ is reported against the best
known $S = 2$ controller \emph{per cell} over a fixed candidate set, not against
the published constants alone. After the seed study of
Section~\ref{sec:results-c} the rule in force is: the baseline in each cell is
the best of the three class-searched rows in that cell, the row achieving it is
named, and the spread between the three is the baseline's uncertainty and is
reported alongside rather than averaged away. The published constants are
retained as the only \emph{enumerated} row and as the reference the baseline is
measured against.

\paragraph{The upper-bound rule.} The anchor's four-constant grid is exhaustive
at its resolution, so that row alone is tight; at eight constants the grid is
$20^8 \approx 2.6\times10^{10}$ points and is not enumerable. Wherever an optimiser
is used, ``no controller found'' means \emph{not found}, not \emph{not possible}.
Two marks carry this through every table and caption:

\begin{quote}
\dg\ searched, not exhaustive --- an upper bound.\\
\ddg\ hand-designed, not searched --- no evidence about the capability itself.
\end{quote}

\noindent A \ddg\ row is weaker evidence than a \dg\ row: it says one guess did
not work. An unmarked row is enumerated and tight.

\section{Methods}
\label{sec:methods}

\subsection{Robot model and reproduction}

The simulator is a purpose-built model of the e-puck as \citet{gauci2014} model
it in Enki: a disk of diameter 7.4\,cm and mass 152\,g, inter-wheel distance
5.1\,cm, wheel velocities independently settable in $[-12.8, +12.8]$\,cm/s,
control cycle 0.1\,s with physics updated ten times per cycle, and a
\emph{zero-width ray} cast from the robot's front returning the first body it
intersects, at infinite range. The arena is unbounded and all metrics are taken
in the swarm-centroid frame, so bulk translation is factored out.

The inter-wheel distance was initially wrong --- 5.3\,cm assumed against 5.1\,cm
correct --- and the error was caught not by inspection but by a test tying the
wheel constants, the body geometry and the kinematics together against the three
derived quantities the source paper publishes. Those reproduce: $R_0 = 14.45$\,cm
to $5\times10^{-5}$\,m, and $\omega_0 = -0.75$\,rad/s and $\omega_1 =
-5.02$\,rad/s to $5\times10^{-3}$\,rad/s.

Integration is exact-arc: constant wheel speeds over a control step trace a
circular arc integrated in closed form, so one period of the state-0 constants
closes to $10^{-9}$\,m and the result is timestep-independent --- a fifty-fold
refinement of the timestep changes no result, and the source paper's substepping
is a no-op here. Contact is positional relaxation, converged to 0.01\% of a body
diameter by 32 passes. Runs are bit-for-bit reproducible from a (seed, run index)
pair, and the environment fields draw from separate seed streams so that changing
$n$ does not reshuffle the terrain.

\subsection{The metric, and the correction the reproduction forced}

Dispersion is the normalised second moment
$u = 2\sum_i \lVert p_i - \bar p\rVert^2 / (n^2 R^2)$, normalised so a perfectly
packed cluster scores about 1, following the convention the source paper reports
\citep{grahamsloane1990}. \emph{The exact
normalising constant has not been checked against the published figures}, so
absolute dispersion values in this paper are provisional; every claim here is a
comparison between cells that share the normalisation, and those are unaffected.

The gate that reproduces the anchor initially failed at small $n$, and stayed
failed across five hypotheses that were all wrong. The cause was not in the
simulator. It was, first, that the check appealed to the published two-robot
proof, which \citet{steinberg2025} have since shown to be unsound; and second,
that the literature's definition of aggregation is \emph{reaching} a connected
configuration, not being in one at $\tau$. For a pair the difference is a factor of
5.9: in the canonical $n = 2$ cell 0.88 of pairs reach a connected configuration
and 0.15 are still in one at $\tau$. Measured against the
criterion the literature actually uses, the reproduction agrees with it --- 88\%
of pairs and 98\% of five-robot swarms reach a single cluster, against the
$\approx$95.8\% reported for this controller at $n = 2$, and every swarm of ten
or more aggregates in every run.

This paper therefore reports two distinct measurements and never conflates them:
\textbf{reach}, the share of runs that ever formed one connected component, a
proportion carrying Wilson intervals; and \textbf{hold}, dispersion at $\tau$,
carrying percentile-bootstrap intervals. Where rows differ on clean ground we
report a \textbf{hold ratio} --- a row's dispersion under a hostile setting
divided by its own dispersion on flat ground --- so that a row which simply
aggregates better cannot masquerade as robust. Reach is a proportion whose median is 1 whenever
the majority succeed, so a median panel draws a flat line while the probability
falls.

\subsection{Terrain model}

Per wheel $w \in \{L, R\}$:
\begin{equation}
v_w \;=\; v_{\mathrm{cmd},w}\, m_w(x,y) \;-\; g_{\mathrm{eff}}\sin\alpha\,
(\hat h \cdot \hat s),
\label{eq:terrain}
\end{equation}
where $m_w(x,y) = \max\!\big(1 + \theta_m f(x,y),\, m_{\min}\big)$ is a traction
multiplier that \emph{each wheel samples at its own contact point}, and the
gravity term depends on the robot's heading $\hat h$ relative to the uphill
direction $\hat s$. Because the wheels generally sample different values, the
curvature changes and the state-0 circle becomes a distorted loop. That is the
point of the model: a scalar field scaling both wheels equally keeps the
instantaneous centre of rotation fixed --- $v' = mv$, $\omega' = m\omega$, so
$dp/d\theta = v/\omega$ is independent of $m$ --- and a constant drift is a
Galilean shift leaving every line-of-sight reading unchanged. Both are held by
tests, and Section~\ref{sec:results-a} measures the difference rather than
arguing for it.

The field $f$ is lattice value noise with quintic interpolation: stateless,
exactly reproducible from a seed, output in $[-1,1]$ but bell-shaped rather than
uniform, with median $|f| = 0.3476$ over 40 seeds and 1.6\,M samples. The
correlation length $\lambda$ is the lattice spacing in metres and features are
roughly that size; there is no closed-form correlation function, and $\lambda$
should be read as a lattice spacing rather than a Gaussian correlation length.
$g_{\mathrm{eff}} = 0.4$ where slope is used; per-wheel slip noise is available
and is zero in every sweep reported here.

Two conventions matter. First, the sign: we keep the minus sign of
Equation~\ref{eq:terrain} and define $\hat s$ as the \emph{uphill} direction, so
heading uphill costs speed and heading downhill gains it. Second, the ceiling:
the field reaches $|f| = 1$, so at $\theta_m = 1$ the multiplier can reach zero,
a wheel stalls and the robot pivots --- a different dynamical regime, not terrain
deformation. All terrain sweeps are cut at $\theta_m = 0.9$, below the
$1/\max|f| = 1.0$ at which that becomes possible, and a floor $m_{\min} = 0.05$
is a safety net verified never to bind inside the swept range. An earlier figure ran to
$\theta_m = 1.5$, where every row degraded six- to fourteen-fold and the ordering
collapsed; that band is stall, and was removed from the swept range rather than
explained.

The $S = 4$ row's extra sensor state is a terrain bit, $|m(x,y) - 1| > \theta_{
\mathrm{bit}}$ at the robot centre, with $\theta_{\mathrm{bit}}$ set to the field
median $\theta_m \times 0.3476$ so that the bit splits roughly 50/50 --- a bit
that is almost always the same value carries nothing, and a capability row built
on one would be measuring something else.

\subsection{Pursuer model}

A single pursuer with omnidirectional perception inside a finite range $r_p$, not
a line-of-sight ray. The asymmetry is deliberate: the robots are the minimal
agents under study and the pursuer is part of the environment. Its dials are the
speed ratio $\rho = 1.5$ throughout (so a locked target cannot escape),
$r_p \in \{0.10, 0.20, 0.35, 0.60, 1.00\}$\,m, confusion $\kappa \in \{0, 0.5,
1.0, 2.5, 5.0\}$ entering as $p_{\mathrm{lock}} = 1/(1 + \kappa
n_{\mathrm{local}})$ with $n_{\mathrm{local}}$ counted inside 0.5\,m, and a
Holling handling time $h$. Targeting is nearest-visible and the search behaviour
when nothing is in range is a correlated random walk; alternatives (fewest
neighbours, random visible, Archimedean spiral) are implemented and were not
swept. Capture is at 0.08\,m centre to centre.

Handling time is required rather than optional, and it is the most consequential
correction the pursuer model needed. With capture free on contact a packed
cluster is a buffet --- one robot per control step, twenty gone in two seconds ---
so aggregation is maximally bad under any metric. The values are set against the
pursuer's own travel time between touching neighbours in a formed cluster,
7.4\,cm at $\rho v_{\max} = 19.2$\,cm/s $= 0.385$\,s: $h = 0.39$\,s is comparable
to it and $h = 1.93$\,s is five times it. The first sweep used $h = 5.0$\,s,
which is why its survival figures are not comparable cell for cell with the
later ones.

\paragraph{Reading $r_p$.} The pursuer's range is meaningless on its own and is
reported in two units throughout: the swarm starts inside a disc of radius
$R = 0.74$\,m at $n = 20$, so the five ranges are 0.14, 0.27, 0.47, 0.81 and
$1.35\,R$. \textbf{$r_p \ge R$ is the perfect-perception corner}: the pursuer sees
the whole starting swarm from anywhere in it, so that column is the degenerate end
of the axis, not the hard one --- the cell where hiding is impossible rather than
the most demanding one.

\subsection{Occlusion}

Two rates on the line-of-sight reading, both optionally spatially correlated by
a smoothed field. It was run once, as a harness shakedown, and contributes one
methodological result to this paper: the modulation is multiplicative and
clamped, so the \emph{realised} rate can run far above the nominal
one --- at a nominal 0.6 its median across runs is \textbf{0.89} and its mean
0.856 --- because the field is
sampled where the robots actually are and a cluster that settles in a bad patch
draws its readings from that patch. Correlated and i.i.d. dropout must therefore
be compared at matched realised rate, and this is a warning about \emph{any}
dial whose intensity varies in space.

\subsection{Controller families and search protocol}

Every searched row used sep-CMA-ES \citep{ros2008}, named precisely: the
separable variant restricts the covariance to its diagonal and is therefore a
\emph{weaker} searcher than full CMA-ES, which is the safe direction --- it
loosens an upper bound and never falsely tightens it. Constants are
box-constrained to $[-1,1]$ with candidates repaired to the box and charged a
quadratic penalty; the objective is the median over a candidate's trials, because
a mean sits between the two modes of a bimodal cell.

Seed discipline is uniform: searches run on training seeds and rows are evaluated
on a \emph{disjoint} seed range, and two rows differing in exactly one respect
share a training seed base so the comparison is about that dial rather than a
different draw of arenas. Single-condition searches used 600 evaluations
$\times$ 12 runs; class searches used 1200 $\times$ 12, doubled and stated,
because giving both the same budget would confound ``no transferring controller
exists'' with ``the class search was under-resourced''.

\paragraph{The class objective.} With one training condition the objective is the
median final dispersion; with several it is the \textbf{geometric mean of the
per-condition medians}. Median dispersion across the class runs from about 1.2 to
tens, so an arithmetic mean --- or a median over pooled runs --- is the hardest
condition wearing a disguise and lets the optimiser abandon the rest of the class,
while averaging logs weights a proportional improvement equally everywhere. It
reduces exactly to the single-condition objective at class size one, so both are
one procedure rather than two incomparable protocols. Neither the geometric-mean
objective nor the doubled budget was ablated: the row that transfers used both,
and this paper cannot say which did the work.

\paragraph{Honest re-scoring.} A search reports the best objective value it saw,
which is the minimum of a noisy sample: the best-looking candidate is partly the
luckiest. We therefore re-score returned controllers at high replication on their
own training seeds and report the gap. Section~\ref{sec:results-c} gives two
worked cases, one small and one large.

\subsection{Statistics and the validity window}

At least 100 runs per cell throughout. Continuous quantities are reported as
medians with 95\% percentile-bootstrap intervals; proportions as Wilson score
intervals. Bootstraps are seeded so figures regenerate identically. Where two
rows share a seed base, comparisons are \emph{paired by run index}: run $i$ is
the same initial placement and the same field in both, so a difference or ratio
is taken run by run rather than between independent medians.

One window has to be declared in advance because it changes what a peak means. A
hold ratio measures ``how much looser is the cluster'' only while there
\emph{is} a cluster; where reach collapses, the ratio is driven by runs that
never aggregated and a maximum in it is a peak in failure. We therefore read a
peak only where \textbf{reach at $\theta_m = 0.9$ is at least 0.8}. The threshold
was fixed after seeing one sweep in which it mattered, so its use \emph{there} is
post hoc and is labelled as such; it was fixed in advance for every sweep
reported after that one.

\section{Results A --- terrain}
\label{sec:results-a}

\subsection{Small perturbations do not help}

The hypothesis this section had to dispose of first is the one carried over from
the discrete model: that a small amount of terrain, like a small amount of motion
noise, should \emph{help} by breaking deadlocks. It does not, and the answer is
the paper's strongest negative result.

It was tested three independent ways: an actuation-noise probe found no
improvement; the coarse terrain grid showed degradation rising monotonically from
zero with no beneficial shoulder; and a dedicated fine sweep of the neighbourhood
of zero --- slope $\alpha \in \{0, 0.5, 1, 1.5, 2, 3\}^{\circ}$ crossed with
$\theta_m \in \{0, 0.02, 0.05, 0.1\}$ at \textbf{200 runs per cell}, twice the
protocol minimum --- put the whole grid between \textbf{1.385 and 1.430} in median
dispersion, with \textbf{0 of 23 non-baseline cells} having an interval disjoint
from the clean baseline's. Every interval overlaps every other; the zero-slope row
reads 1.406 \CI{1.389}{1.429}, 1.388 \CI{1.360}{1.406}, 1.385 \CI{1.371}{1.414},
1.414 \CI{1.399}{1.436}.

Three nulls at that replication are a bound on the effect's size, not a failure to
find it. The analogy fails because the published mechanism is noise perturbing a
\emph{precise balancing of forces} so robots can push past one another --- a
contact deadlock in a discrete model --- and this continuous model shows no sign
of such a deadlock at $n = 20$. Whether the deadlock exists here at all was never
tested directly, which is the one thing this result does not settle.

\subsection{The capability minimum does not move}

Over $0 \le \theta_m \le 0.9$ at $\lambda/R_0 = 0.69$, $c = (2,0,0,0)$ suffices
at every point, at a start radius of 0.74\,m and at 1.5\,m. Three measurements
establish it, and the third is constructive.

First, the extra sensor state does not buy back the degradation. At the worst
cell the searched $S = 4$ row sits at a hold ratio of 1.212 \CI{1.158}{1.251}
against the searched $S = 2$ row's 1.099 \CI{1.044}{1.145} --- \emph{disjoint}
intervals,
with the cheaper row ahead (Figure~\ref{fig:h3}). The honest statement is not that
$S = 4$ is worse, since both rows are upper bounds and 600 evaluations over eight
constants is a thinner search per dimension than 600 over four; it is that
\textbf{no $S = 4$ controller better than the searched $S = 2$ one was found at
equal budget}.

Second, warm starting removes the direction of that objection. A search started
at an eight-constant table whose two halves are equal --- an $S = 4$ controller
that ignores its own terrain bit and behaves exactly like the $S = 2$ optimum ---
begins holding that optimum, so the only question it can be asked is whether
\emph{using} the bit beats \emph{ignoring} it. The search did move (L2 distance
0.464 from its start, halves diverging by up to 0.254, so the returned controller
genuinely uses the bit) and it ties: absolute dispersion at the worst cell 1.412
\CI{1.349}{1.503} against 1.383 \CI{1.350}{1.441}, overlapping, with the point
estimate 2.1\% \emph{worse}. Its reported training objective, 1.2595, was better
than either cold search's (1.3012 and 1.3009) and bought nothing on held-out
seeds --- the first of three encounters in this paper with a search's own score
as an optimistic estimate.

Third, a hand-built composite that switches between the two $S = 2$ rows on the
terrain bit is worse than either at every $\theta_m$ --- 1.999 \CI{1.861}{2.143}
against 1.383 \CI{1.350}{1.441} at $\theta_m = 0.9$ --- because half the time it
deliberately selects the worse of its two behaviours. Its $\theta_m = 0$ value
reproduces the published constants exactly (1.427 \CI{1.399}{1.467}), which is
the internal consistency check: on flat ground the bit is never set, so the
composite \emph{is} the anchor.

These three are quoted in two different units and the switch matters: the first
pair are hold ratios --- throughout this paper, the median of the per-run ratio
of a row's dispersion to its own flat-ground dispersion, paired by run index ---
while the second and third are \emph{absolute} dispersion at $\theta_m = 0.9$. A
hold ratio of 1.099 and an absolute dispersion of 1.383 describe the same row.

\begin{figure}[t]
\centering
\includegraphics[width=\textwidth]{terrain_h3_capability.png}
\caption{The terrain-sensing state does not buy back the peak; re-searching four
constants does. Rows: the enumerated anchor, a searched $S = 2$ row\,\dg\ and a
searched $S = 4$ row\,\dg\ with a terrain bit, at equal budget and on held-out
seeds. Reach and hold are plotted separately because they answer different
questions.}
\label{fig:h3}
\end{figure}

\subsection{Terrain is a performance tax, and its size is a property of the
controller}

What moves along this dial is the parameter setting inside a fixed capability,
and most of what a re-search buys has nothing to do with terrain. The
decisive control is a second searched row, identical in optimiser, budget and
training seed base, differing only in the $\theta_m$ of its training condition:
one row trained on rough ground and one trained on flat ground that never saw
terrain at all.

At $\theta_m = 0.9$, hold ratios against each row's own flat baseline are 2.149
\CI{1.853}{2.455} for the anchor, 1.200 \CI{1.124}{1.272} for the flat-trained
row\,\dg, and 1.099 \CI{1.044}{1.145} for the rough-trained row\,\dg. \textbf{A
controller that never saw terrain degrades by 20\% where the published constants
degrade by 115\%.} The anchor's hold ratio therefore overstates what terrain does
to a well-matched controller by $1.96\times$.

Decomposing the gap between the anchor and the rough-trained row at
$\theta_m = 0.9$ --- 3.150, 1.438, 1.383 in absolute dispersion --- gives
\textbf{96.9\% objective-tuning and 3.1\% terrain-tuning at a 0.74\,m start
radius, against 99.7\% and 0.3\% at 1.5\,m}. The two figures travel together in
every statement of this result, because the terrain share is the part that does
not survive a change of initial condition. This is a decomposition of
\emph{levels}, and it decomposes only because medians of levels are additive; it
is not a hold ratio. What the paired data adds is a test on the terrain term
itself, and that term is \textbf{not distinguishable from zero at either start
radius}: $+0.0818$ \CI{-0.0008}{0.1331} at 0.74\,m and $-0.0124$
\CI{-0.1116}{0.0775} at 1.5\,m. On flat ground it reverses sign and costs
4.1\%. So terrain-specific tuning is worth about $\pm 4\%$ of dispersion in
either direction and cannot be shown to differ from nothing, while matching the
controller to the objective at all is worth 119\%.

\begin{figure}[t]
\centering
\includegraphics[width=\textwidth]{terrain_tuning_control.png}
\caption{The $2\times2$ tuning control. The two searched rows\,\dg\ differ only
in the $\theta_m$ they were trained at --- same optimiser, budget and training
seeds. The bottom panel is the difference taken run by run, which is available
because the two rows share a seed base.}
\label{fig:tuning}
\end{figure}

\subsection{The mechanism is per-wheel curvature, not speed heterogeneity}

A scalar speed field still makes robots in different places move at different
rates, and that alone changes who meets whom. A paired control settles whether
that is the mechanism: both rows use the same placements and the same traction
field and differ only in whether each wheel samples at its own contact point or
both take one value sampled at the robot's centre. The rule was fixed before
running: if the scalar row reproduces the per-wheel row's peak location
\emph{and} magnitude within interval, the curvature explanation is retired.

It does not reproduce. Pooled over every $\theta_m > 0$ cell (1800 runs each),
the per-wheel row degrades by 1.175 \CI{1.154}{1.190} and the scalar row by
0.980 \CI{0.976}{0.986}. The per-wheel row peaks at 2.21 \CI{1.90}{2.41} at
$\lambda/R_0 = 0.69$; the scalar row has \emph{no peak at all} and is flat within
$[0.96, 1.01]$ across the entire grid (Figure~\ref{fig:mech}). The curvature
explanation stands, and the model correction that motivated it is now supported
by a control rather than by an argument.

One thing is reported without explanation, as the rule required: the scalar row
sits slightly \emph{below} 1.0, an interval excluding 1.0, so a pure scalar speed
field appears to aggregate about 2\% better than flat ground. That is not in the
direction the noise-helps hypothesis predicted, 2\% is at the edge of what this
design resolves, and it remains unexplained.

\begin{figure}[t]
\centering
\includegraphics[width=0.92\textwidth]{terrain_mechanism.png}
\caption{Paired mechanism control: identical placements and identical field,
differing only in how the field is sampled. The scalar-centre row is the
rejected model; it produces no degradation anywhere.}
\label{fig:mech}
\end{figure}

\subsection{The mechanism lemma: what the residual is made of}

With zero slope the per-timestep heading-rate residual is exactly
\begin{equation}
\mathrm{residual} \;=\; \omega_{\mathrm{actual}} - \omega_{\mathrm{commanded}}
\;=\; v\,\frac{m_R - m_L}{\ell} \;+\; \frac{\omega}{2}\,(m_R + m_L - 2),
\end{equation}
and to first order about the robot centre, with $\hat n$ running from the left
wheel contact to the right,
\begin{equation}
\mathrm{residual} \;=\; \underbrace{v\,\partial m/\partial n}_{\text{gradient}}
\;+\; \underbrace{\omega\,(m(\text{centre}) - 1)}_{\text{mean traction}}.
\label{eq:decomp}
\end{equation}
Both terms are the \emph{same} per-wheel mechanism --- the wheels seeing
different ground, and a turning robot on ground uniformly slower than nominal ---
and this is arithmetic verified to six decimals against the simulator, not a
second hypothesis. Because the two terms are built from the same field they are
correlated, so regressing on the gradient alone yields a \emph{biased} slope
rather than a merely noisy one, and both fits are reported.
$\partial m/\partial n$ is differenced from the field at the robot centre, never
from the wheel samples, since regressing on those would be an identity.

Regressed over twelve million robot-timesteps per point, the full expansion has
slope \textbf{0.9943} with $R^2$ \textbf{0.9999} at $\ell/\lambda = 0.128$: the
residual \emph{is} the first-order per-wheel traction effect and there is
essentially nothing else in it. Slope falls monotonically with $\ell/\lambda$ ---
0.9762, 0.90, 0.62 --- and collapses to \textbf{0.1351} with $R^2$
\textbf{0.0870} at $\ell/\lambda = 2.04$, where a first-order expansion about the
centre has stopped describing what the wheels see. The gradient-only fit gives
slopes 0.52--0.81 and $R^2$ 0.02--0.28 at every $\ell/\lambda$, which is the
predicted consequence of omitting a correlated term roughly ten times its size.

The dependence is on the \emph{ratio}, not on the axle. Crossing three axle
lengths with three correlation lengths so that the same $\ell/\lambda$ is reached
by different $(\ell, \lambda)$ pairs, the slope at matched $\ell/\lambda$ agrees
to \textbf{0.0201} over a fourfold range of $\lambda$ --- 2.3\% of the 0.86 range
the trend covers --- and to 0.0001 at $\ell/\lambda = 0.255$
(Figure~\ref{fig:collapse}).

\begin{figure}[t]
\centering
\includegraphics[width=\textwidth]{terrain_lambda_collapse.png}
\caption{The linearisation collapses on $\ell/\lambda$ and not on the axle: the
same ratio reached by different (axle, $\lambda$) pairs gives the same slope to
0.0201. Four of the nine points come from hand-picked geometry rows\,\ddg, and
the widest-axle row places the wheel contacts outside the body.}
\label{fig:collapse}
\end{figure}

\subsection{The worst correlation length is not the turning circle}

The hypothesis this project began with was that terrain bites hardest when the
correlation length is comparable to the controller's turning circle, at
$\lambda/R_0 \approx 0.7$. The first measurement appeared to find exactly that, and
it takes two separate statements to say what it actually found, because the peak
has a \emph{magnitude} and a \emph{location} and $R_0$ governs only one of them.

\textbf{Magnitude scales with $R_0$, and not with the axle.} At $\theta_m = 1.0$,
varying the axle fourfold at fixed $R_0$ moves peak degradation by 5\% (2.97,
2.83, 2.93); varying $R_0$ fourfold at fixed axle moves it by 560\% (2.97, 1.59,
10.46). How bad the worst case gets is a property of the turning circle.

\textbf{Location does not.} The apparent ratio law --- a peak at
$\lambda/R_0 \approx 0.69$ for four of five rows --- holds at $\theta_m = 0.7$ and
at $\theta_m = 1.0$ only two of the same five rows sit there. Two measurements
below dispose of it outright.

Two controllers differing threefold in $R_0$ peak at the \emph{same} $\lambda$ in
metres, 7.46\,cm, which is $\lambda/R_0 = 0.52$ for one and 1.57 for the other;
at the second row's predicted worst $\lambda$ its hold ratio is 1.05--1.09, in
the bottom half of its range. Inside a single controller family --- three rows of
the anchor's table differing \emph{only} in the state-0 forward constant, so that
$R_0$ is the only thing that moves --- doubling $R_0$ moves the peak by
\textbf{1.39$\times$} \CI{0.72}{1.93}, an interval that \emph{excludes the ratio
law's 2.00} and contains a fixed scale's 1.00. Proportionality is rejected
(Figure~\ref{fig:r0family}). A weak sub-proportional dependence is not excluded:
the point estimate lies between 1 and 2, and peaks are located only to a log grid
of ratio 1.39.

\begin{figure}[t]
\centering
\includegraphics[width=\textwidth]{terrain_lambda_r0_family.png}
\caption{Peak correlation length within one controller family, at two start
radii. The three rows differ only in the state-0 forward constant, so $R_0$ is
the only length that moves. The shaded band is the pre-declared validity window
(reach $\ge 0.8$ at $\theta_m = 0.9$); its use at the larger start radius is post
hoc and the decision is read from the companion sweep at the smaller one. All
three rows are numerical devices\,\ddg\ built from the anchor's table.}
\label{fig:r0family}
\end{figure}

What the length \emph{is} appears to be the robot. With $R_0$, the axle, the
sensor model and the start radius all held, a 7.4\,cm body degrades worst at
$\lambda = 7.46$\,cm and a 14.8\,cm body at 14.39\,cm --- \textbf{1.01 and 0.97
body diameters}, log-log slope \textbf{0.948} \CI{0.474}{1.897}, excluding 0 and
containing 1 (Figure~\ref{fig:body}). This result is reported at SUGGESTED
strength and the hedge is part of the claim: it rests on the two body sizes that
aggregate at this start radius, the third reaches a cluster in at most 72\% of
runs and has a flat-ground dispersion of 3.057 against 1.427 for the base body,
and including it anyway lowers the slope to 0.474. Three lengths scale with the
body and this experiment separates none of them --- the body as the size of an
obstacle, the cluster link distance at 1.5 body diameters, which is the
\emph{metric's} own definition of ``together'', and the occlusion footprint,
since a larger body blocks more line of sight. No mechanism is proposed.

\begin{figure}[t]
\centering
\includegraphics[width=\textwidth]{terrain_lambda_body.png}
\caption{Peak $\lambda$ against body diameter, with $R_0$, the axle, the sensor
model and the start radius held. The smallest body does not aggregate at this
start radius and is excluded by the pre-declared validity window; hollow markers
are outside it. Both non-default body sizes are numerical devices\,\ddg\ rather
than buildable robots.}
\label{fig:body}
\end{figure}

\section{Results B --- the pursuer}
\label{sec:results-b}

\subsection{Confusion acts through aggregation, and a control proves it}

The first pursuer sweep found survival rising steeply with confusion and read it
as dilution. But every row in it aggregated identically when no pursuer was in
view, so nothing separated \emph{the swarm aggregated} from \emph{the pursuer's
perception is poor}. The missing control is a row with byte-identical pursuer
sensing and no clustering: a 99\,cm state-0 arc instead of the anchor's 14.45\,cm
circle, so a robot never stops for a neighbour and no cluster forms, with a flee
response copied verbatim from the aggregating row it is compared against. It
differs from that row in exactly one respect.

As $\kappa$ goes from 0 to 5, mean survival multiplies by \textbf{3.05} and 2.23
for the two aggregating rows at $h = 0.39$\,s, and by \textbf{4.10} and 1.91 at
$h = 1.93$\,s. The dispersive control gains \textbf{1.20} and \textbf{1.13}
(Figure~\ref{fig:kappa}). Confusion enters the model only through
$p_{\mathrm{lock}} = 1/(1 + \kappa n_{\mathrm{local}})$ and only a clustered
swarm has a large $n_{\mathrm{local}}$, so this is the mechanism doing what it is
written to do --- but it had not been \emph{shown} before, because nothing
separated it from the pursuer simply being worse at its job.

Handling time sharpens it in the predicted direction: raising $h$ from comparable
to the pursuer's inter-neighbour travel time to five times it grows the blind
row's $\kappa$-response from $\times3.05$ to $\times4.10$ while the dispersive
row's shrinks. Dilution needs the predator busy, and making it busier makes
clustering pay more.

\begin{figure}[t]
\centering
\includegraphics[width=0.92\textwidth]{pursuer_dispersive_kappa.png}
\caption{Confusion helps the aggregating rows two- to fourfold and the
dispersive control 13--20\%. The two hand-designed rows\,\ddg\ have identical
pursuer sensing and differ only in spatial strategy.}
\label{fig:kappa}
\end{figure}

\subsection{Where aggregation is a liability, and where it dominates}

Survival here is the \textbf{mean per-robot} survival with a Wilson interval on
the pooled robot count, not a median of per-run fractions: at $n = 20$ a per-run
fraction can only land on a twentieth, so its interval reports that grid rather
than the uncertainty.

Two comparisons are available and only one is about aggregation. The blind row
falls to 0.3636 and then to 0.0199 as the pursuer's range grows, while the
dispersive row holds 0.6238, 0.5084 and 0.3816 --- but the blind row has neither
pursuer sensing nor a flee response, so that pair varies capability and spatial
strategy at once.
\textbf{The matched comparison is the ternary row\,\ddg\ against the dispersive
row\,\ddg}: identical pursuer sensing, same flee response, differing only in
whether the swarm clusters. It is what the control was built for and what the
claim must rest on.

\begin{table}[h]
\centering
\caption{Mean per-robot survival pooled over $\kappa$ at $h = 1.93$\,s, with
Wilson intervals on 10\,000 robots per cell. The first column varies capability
and strategy together; the last two differ only in spatial strategy and are the
comparison that speaks to aggregation. Both hand-designed rows are \ddg.}
\small
\begin{tabular}{llll}
\toprule
$r_p$ ($r_p/R$) & blind & ternary\,\ddg & dispersive\,\ddg \\
\midrule
0.10\,m (0.14) & 0.8077 \CI{0.7999}{0.8153} & 0.8788 \CI{0.8723}{0.8851} & 0.8621 \CI{0.8552}{0.8687} \\
0.20\,m (0.27) & 0.6164 \CI{0.6068}{0.6259} & 0.7369 \CI{0.7282}{0.7454} & 0.7442 \CI{0.7356}{0.7527} \\
0.35\,m (0.47) & 0.3636 \CI{0.3542}{0.3731} & 0.4875 \CI{0.4777}{0.4973} & 0.6238 \CI{0.6143}{0.6332} \\
0.60\,m (0.81) & 0.0199 \CI{0.0173}{0.0228} & 0.1295 \CI{0.1231}{0.1362} & 0.5084 \CI{0.4986}{0.5182} \\
1.00\,m (1.35)\rlap{\,\,perfect perception} & 0.0199 \CI{0.0173}{0.0228} & 0.0655 \CI{0.0608}{0.0705} & 0.3816 \CI{0.3721}{0.3912} \\
\bottomrule
\end{tabular}
\end{table}

On the matched pair, aggregating wins at $0.14\,R$ with disjoint intervals
(0.8788 against 0.8621), the two overlap at $0.27\,R$ (0.7369 against 0.7442),
and dispersing wins from $0.47\,R$ outward with disjoint intervals. \textbf{The
crossover for the matched pair lies between $0.27\,R$ and $0.47\,R$}, and the
paper cannot place it more finely: a denser range axis was never run. Confusion
narrows the remaining gap --- mean survival at $\kappa = 5$ is 0.592, 0.628 and
0.666 across the three rows --- but does not close it. That aggregation loses on
average \emph{over this grid} is a fact about the grid as much as about
aggregation: two of five columns are at or past half the start radius and one is
the degenerate corner, so the upper half is deliberately brutal, and the
qualification belongs with the claim.

Survival alone is the wrong axis. Dispersion among survivors for the dispersive
row is \textbf{380--505} against \textbf{1.43--1.63} for every aggregating row:
the price of its survival, in the units of the task the swarm exists to perform.
It is the task-abandoning end of the front, and reading it as ``the best row''
treats the axis it gives up as free. \textbf{At $r_p = 0.35$\,m
$= 0.47\,R$ \emph{and} $\kappa = 3$ there is no trade-off at all}: every
aggregating row beats it on \emph{both} coordinates --- survival 0.7565--0.7765
against 0.6510, all four disjoint, and dispersion about 1.5 against 386 --- so the dispersive row is
Pareto-dominated (Figure~\ref{fig:pareto}).

\textbf{The $\kappa$ in that sentence is doing the work, not the range.} It is
the same $0.47\,R$ column in which the blind row is at 0.3636 and the ternary row
at 0.4875 once survival is pooled over $\kappa$; at $\kappa = 3$ alone all four
aggregating rows sit at 0.7565--0.7765, every one of them disjointly above the
dispersive row's 0.6510. That is exactly what the control predicts: the
dispersive row is nearly $\kappa$-insensitive, gaining 1.20 and 1.13 from
confusion, while the aggregating rows swing 3.05 to 4.10. \textbf{The ordering at
a given pursuer range is therefore a function of $\kappa$, not of range alone},
and no statement of the form ``aggregation wins inside this radius'' is
well formed without one. Where confusion is strong, confusion has made clustering
safe and clustering is also what the task wants.

Where survival saturates, timing still separates the rows: at $r_p = 0.2$\,m,
$\kappa = 0$ the median time to wipeout runs 47.1, 63.8, 81.0 and 109.8\,s across
the four aggregating rows, recovering at pinned cells the ordering the survival
figures show elsewhere. Saturation is stated rather than averaged over: 29\% of
runs in this sweep ended in a wipeout, concentrated exactly where the claim is
strongest.

\begin{figure}[t]
\centering
\includegraphics[width=\textwidth]{pursuer_pareto.png}
\caption{Survival against base-task performance among survivors, at three cells.
$r_p$ is given in metres and in start radii; the third cell is the
perfect-perception corner. Four of five rows are hand-designed\,\ddg, so the
front is between spatial strategies rather than between capabilities.}
\label{fig:pareto}
\end{figure}

\subsection{What one sensor state buys, and what that claim can bear}

Among the rows tried, going from $S = 2$ to $S = 3$ --- being able to tell a
pursuer from a robot --- is worth more than anything else on the capability axis.
The numbers here come from the \emph{first} pursuer sweep, which used a long
handling time of $h = 5.0$\,s and which is pooled differently from the survival
table above: a single median over the whole $r_p \times \kappa$ grid, rather than
one figure per range column. Mean per-robot survival over that grid is 0.4324
\CI{0.4281}{0.4367} for the blind row, 0.5851 \CI{0.5807}{0.5894} for the ternary
row\,\ddg\ and 0.6079 \CI{0.6036}{0.6122} with a memory bit added\,\ddg, on
50\,000 robots each. The blind row's 0.4324 here and its 0.0199 in the two
longest-range columns of the table above are the same row under a different
handling time and a different pooling, not a disagreement. The blind row is wiped out at
$\kappa = 0$ for every range beyond 0.10\,m while any row that can see the
pursuer keeps three to eighteen robots.

This is stated at SUGGESTED strength with its hedge attached, because the rows
above the blind one are hand-designed\,\ddg\ rather than searched: a lower bound
on what $S = 3$ offers, not a measurement of it. The same caution kills a tempting
reading in the other direction: a five-state row scores \emph{below} the
three-state row over most of the grid (0.5188 \CI{0.5144}{0.5232} against 0.5851
\CI{0.5807}{0.5894}), and
that is a hand-design artefact. Its table is a written guess and a bad one; it
says no good five-state controller was found by hand, and reporting it as ``more
sensing hurts'' would be the single easiest mistake to make with this data.

Finally, a coverage statement rather than a result: \textbf{no $K \ge 1$ row was
run}. Communication is not wired in the simulator --- the receive channel is held
at zero --- so a $K = 1$ row would silently behave as $K = 0$, and a row that
quietly does nothing is worse than an absent one. This paper covers $S$ and $M$,
not $K$.

\section{Results C --- methodology}
\label{sec:results-c}

This section shares no apparatus with the two above beyond the simulator and the
anchor controller. It is about how a searched controller is obtained and reported
rather than about what any environment does to a swarm, and it can be read on its
own.

\subsection{A controller searched at one operating point need not work at
another}

The two rows searched at $n = 20$ and a 0.74\,m start radius beat the published
constants comfortably \emph{there}. Moved, they fail. At four times the training
start radius, under terrain, with the trial length extended sixfold so that
slowness cannot be mistaken for failure, the rough-trained row forms a cluster in
\textbf{34\%} of runs against the anchor's \textbf{100\%}, and its paired
dispersion ratio does not move at all across a sixfold change in trial length
(0.137, 0.141, 0.148). The flat-trained row's collapse in that cell is largely
truncation --- given six times the budget it draws level, 1.297
\CI{0.960}{1.495}, an interval still containing 1 --- but level is not the 1.94
it manages at the tuned radius.

At $n = 50$ the failure is not about trial length at all. On \emph{flat ground},
at every start radius, both searched rows are worse than the published constants:
paired ratios 0.973, 0.870, 0.775 and 0.944, 0.818, 0.777, every interval below
1. Six times the trial length moves the worst of these only from 0.775 to 0.839
\CI{0.795}{0.880}, while the anchor's own dispersion moves 1.204 to 1.209 --- it
was converged all along (Figure~\ref{fig:regime}).

\begin{figure}[t]
\centering
\includegraphics[width=\textwidth]{terrain_regime_robustness.png}
\caption{Reach, dispersion and time to first cluster across start radius, swarm
size and terrain, for the enumerated anchor and two single-condition searched
rows\,\dg. The vertical marks the start radius the searched rows were tuned at.}
\label{fig:regime}
\end{figure}

\subsection{The trade is gathering rate against holding quality}

The explanation is visible on flat ground, where nothing fails. Median time to
first cluster at a 3.0\,m start radius, $n = 20$, is \textbf{140\,s} for the
anchor ($R_0 = 14.45$\,cm), \textbf{245\,s} ($R_0 = 7.47$\,cm), \textbf{280\,s}
(5.53\,cm), \textbf{330\,s} (4.74\,cm) and \textbf{340\,s} (4.47\,cm) ---
monotone in $R_0$ across five independently obtained controllers. The blind state
\emph{is} the search behaviour, so a smaller turning circle holds a formed
cluster more tightly and covers ground more slowly. At the tuned start radius the
cost is invisible: 30\,s against 20\,s inside a 600\,s budget. At four times that
radius it is a factor of two to two and a half, and terrain multiplies travel
time on top of it.

This is an association across five rows, not a controlled measurement: $R_0$ was
never swept as a dial in its own right, and the paper says so. What it does do is
predict, rather than accommodate, the pattern in the next subsection.

\subsection{Searching the mission class produces a controller that transfers}

The failure above is about two controllers, not about the capability: nothing had
asked the optimiser to work anywhere but the arena it was given. Asking it --- the
same four constants, the same optimiser, a doubled and stated budget, and an
objective averaged over the six conditions the evaluation grid is drawn from ---
produces a row that transfers.

Cell by cell on the held-out grid at $\tau = 600$\,s, the class-searched
row\,\dg\ wins 9 of 12 cells, loses 2 and ties 1, against 7/4/1 and 7/5/0 for the
two single-condition rows. It \emph{completely} fixes the swarm-size failure: on
flat ground at $n = 50$, where the single-condition rows score 0.973, 0.870,
0.775 and 0.944, 0.818, 0.777, it scores \textbf{1.026} \CI{1.020}{1.033},
\textbf{1.016} \CI{1.009}{1.023} and 0.999 \CI{0.991}{1.008} --- two wins and a
tie, with absolute dispersion at the hardest of those cells 1.200 against the
anchor's 1.204 where the single-condition row sits at 1.545.

Both losses are at the largest start radius under terrain, a column needing a
longer trial for \emph{every} row: at $\tau = 600$\,s the anchor itself reaches a
cluster in only 46\% of runs there. At $\tau = 3600$\,s the class row wins that
cell at $n = 20$ --- \textbf{1.436} \CI{1.249}{1.673} --- and ties at $n = 50$,
making it \textbf{10 W, 2 tied, 0 L} on dispersion across the twelve cells
(Figure~\ref{fig:class}).

Reach is where the anchor still wins, and it is why the claim carries an
exception. At $\theta_m = 0.9$ with the 3.0\,m cells at $\tau = 3600$\,s, the
class row is \emph{better} on reach at the tuned radius (0.98 \CI{0.93}{0.99}
against 0.86 \CI{0.78}{0.91}), tied at 1.5\,m and at both $n = 50$ cells below
3.0\,m, and \emph{worse} at both 3.0\,m cells: 0.86 against 1.00 at $n = 20$, and
0.95 against 1.00 at $n = 50$. So the class-searched row is the better baseline
everywhere except the largest start radius under terrain, where the enumerated
row still owns the corner.

The mechanism reading predicts which cells those are: training on the class moved
$R_0$ from 5.53\,cm to 7.47\,cm, 35\% of the way back toward the anchor's
14.45\,cm, buying back gathering rate --- and every loss in this section is at the
one cell where reach binds.

\begin{figure}[t]
\centering
\includegraphics[width=\textwidth]{terrain_class_search.png}
\caption{The class-searched row\,\dg\ against the enumerated anchor and the two
single-condition searched rows\,\dg\ on the held-out grid. At $\tau = 600$\,s the
3.0\,m column is truncated for every row, which is why that column is decided on
a separate trial-length series.}
\label{fig:class}
\end{figure}

\subsection{The noise floor: when a search cannot see its own gradient}

A companion experiment gives the sharpest methodological result in the paper, and
it began as a wrong diagnosis. The rough-trained class row did not move $R_0$ at
all, and the first explanation offered was that its training objective was
truncated --- at the training trial length almost no candidate aggregates at the
hardest conditions, so those conditions carry a near-constant penalty with no
gradient in it. Re-running the identical search with the trial length raised at
exactly those conditions, everything else held including the training seed base,
moved $R_0$ from 4.47\,cm to \textbf{4.44\,cm}: not at all. At the cell the change
targeted, training \emph{at} the longer trial length made the row \emph{worse}
(reach 0.26 \CI{0.18}{0.35} against 0.36 for the row trained at the short one).

The real cause is measurable. Scoring four controllers on that search's own
objective, on its own training seeds, at 100 runs per condition instead of 2, the
flat-trained class row scores \textbf{1.6456} against the returned row's
\textbf{2.6553}: a strictly better point on the objective being minimised already
existed, and was found by a different search. The objective does not prefer the
returned controller; the search did not find its own optimum. \textbf{It does
not locate that optimum either}, so ``the objective prefers a large $R_0$'' is
not established by this --- only that the returned row is not its best point.

It failed because of one cell. There the dispersion runs from 1.14 to 242.50 over
100 runs, and the search saw the median of \emph{two} of them per evaluation. The
5th--95th percentile of that two-run estimate spans \textbf{2.79 log units}
(6.98 to 114.07) against a \textbf{0.48 log unit} gap between the returned row and
the better one that existed --- noise \textbf{5.8$\times$} the signal. Under that
noise a best-so-far selection returns whichever candidate drew a lucky pair,
which is exactly why the search's reported best objective, \textbf{1.6009}, sits
\emph{below} an honest re-score of the same constants, \textbf{2.6553}. Raising
the trial length made it worse rather than better: it turned uniform failure into
outcomes spanning two orders of magnitude, which is more gradient \emph{and} far
more variance, and the variance won (Figure~\ref{fig:noise}).

The rule this yields is not ``check for truncation'' but \textbf{check whether
any single condition's per-evaluation estimate is noisier than the differences
the search must resolve}. The remedy that follows is replication per condition,
or a variance-stabilising statistic, rather than a longer trial --- and it
\emph{was not run}, so the remedy follows from the measurement rather than from
a demonstration. Its corollary is that
a training objective must be re-scored at high replication before it is quoted,
and that a searched row's standing must rest on held-out evaluation --- which is
what caught this failure.

\begin{figure}[t]
\centering
\includegraphics[width=\textwidth]{terrain_class_tau.png}
\caption{Why a class search can fail to find its own optimum: the objective
re-scored at high replication (left), what the search returned (middle), and the
two-run estimate the search actually saw at the worst condition against the
signal it had to resolve (right).}
\label{fig:noise}
\end{figure}

\subsection{The baseline is a reproduced region, not a reproduced optimum}

The class-searched row was found by one optimiser seed. Repeating the search with
two further seeds, identical in every other respect including the training seed
base, returns controllers whose turning circles are \textbf{7.47, 8.88 and
8.64\,cm} --- a spread of 1.41\,cm, \textbf{16.9\% of the mean}, outside the
10\% agreement rule declared in advance --- and whose held-out dispersion
intervals are \textbf{disjoint in 4 of 12 cells} (Figure~\ref{fig:seeds}).

The substance survives replication and is strengthened by it. All three seeds
beat the enumerated anchor in at least 9 of 12 cells (9/2/1, 10/0/2, 10/1/1
W/L/tied), all three fix the flat-ground failure at $n = 50$ that defeated every
single-condition row, and all three move $R_0$ up toward the anchor's, so the
\emph{direction} of the gather-versus-hold trade reproduces in every draw. What
does not survive is naming one row as \emph{the} baseline: the row originally
named is the weakest of the three by cell count, and no seed dominates another
head to head. The baseline is therefore best-of-three\,\dg\ per cell, and the
spread is its uncertainty.

Three of the four disjoint cells are flat ground at $n = 50$, where differences
of 0.02--0.03 in dispersion are resolvable at 100 runs and practically small. The
fourth is the truncated terrain cell, where the spread is threefold at
$\tau = 600$\,s (12.517, 3.701, 13.701) and closes to 0.011 in the paired ratio
at $\tau = 3600$\,s (1.436, 1.430, 1.441, all three beating the anchor). The
paper reports it at both trial lengths, because a reader who saw only one would
either overstate or understate the scatter.

The seed study also bounds the winner's curse where it is small. Each search's
reported objective is optimistic by \textbf{+0.019, +0.032 and +0.018} --- 1.5\%
to 2.7\% --- against the \textbf{+1.05} measured for the rough class above. That
is the noise-floor account predicting its own boundary: on flat ground every
condition reaches a cluster in every run and a two-run median is a usable
estimate; at high terrain amplitude it is not.

\begin{figure}[t]
\centering
\includegraphics[width=\textwidth]{phase0_seed_reproducibility.png}
\caption{Three searches differing only in the optimiser seed. Constants returned
(left), behaviour on the held-out grid at both swarm sizes, and the gap between
each search's reported objective and an honest re-score of the same constants.}
\label{fig:seeds}
\end{figure}


\section{Discussion}
\label{sec:discussion}

\subsection{Two kinds of hostility, two kinds of answer}

The two environment axes behave differently, and that difference is this paper's
organising result. \textbf{Passive hostility taxes the minimal
controller without moving it; active hostility moves the minimum along the dial
the spatial strategy can act on.} Terrain --- an environment that resists but does
not respond --- degrades aggregation quality by a factor set almost entirely by
the controller's match to its objective, and never makes two sensor states
insufficient or four sufficient-and-better. A pursuer --- an environment that
responds --- changes which spatial arrangement survives, and the one capability
separating the rows we could build by hand is the sensor state that tells a
pursuer from a robot.

The first half rests on three measurements with disjoint intervals: the two-state
controller suffices at every point of the terrain dial at both start radii
tested; the four-state search never found a better controller at equal budget,
including from a warm start at the two-state optimum; and the composite that
switches between a flat-tuned and a rough-tuned table on a terrain bit is worse
than either of its own halves everywhere. Three independent routes to the same
negative answer is the strongest form this claim can take without an exhaustive
search of the four-state space, which this dimensionality does not permit.

The second half is weaker by construction. Every pursuer row above the baseline
is hand-designed, so the capability statement is a \emph{lower bound on what three
sensor states offer, not a measurement of it}, and the five-state row's poor
showing is a hand-design artefact rather than evidence that more sensing hurts.
What is not weak is the mechanism: the dispersive control shares the pursuer's
sensing byte for byte and gains $1.20$ and $1.13$ from confusion where the
aggregating rows gain $3.05$ and $4.10$. Confusion is a property of the group the
robots form, not of the sensor they carry.

\subsection{The tax is mostly not about terrain}

The most transferable finding for practitioners is the least intuitive one.
Nearly all of what looks like ``the controller was re-tuned for terrain'' is the
controller being re-tuned for \emph{this dispersion objective at this start
radius}, which the published constants were never tuned for --- they were found by
exhaustive search on a different criterion. The decomposition is $96.9\%$
objective-tuning against $3.1\%$ terrain-tuning at $0.74$\,m, and $99.7\%$ against
$0.3\%$ at $1.5$\,m. The terrain-specific share is small at the tuned start
radius, vanishing at twice it, and on flat ground it reverses sign: there the
rough-trained row \emph{costs} 4.1\% against the flat-trained one, so the
flat-trained row is the better flat-ground controller of the two. Terrain-specific
tuning is worth about $\pm 4\%$ of dispersion in either direction; matching the
controller to the objective at all is worth 119\%. A reader who takes ``re-tuning solves terrain'' from the raw hold ratios
takes the wrong lesson; the right one is that terrain-specific tuning buys very
little, and that most of the apparent benefit is available without ever
simulating terrain.

A terrain robustness study reporting only a re-tuned controller against a
published baseline will attribute to the environment an effect that is
overwhelmingly objective mismatch. Once the objective is matched, terrain costs
\textbf{10--20\%} in held dispersion --- 1.099 for the rough-trained row and
1.200 for the flat-trained one --- a performance tax a designer can budget for,
not a capability requirement.

\subsection{Implications for field minimal swarms}

Three practical statements follow, at the grades the register gives them.

\emph{Design over the mission class, not over one arena}, which is the cheapest
of the three to adopt: a geometric mean of per-condition medians over six
conditions spanning start radius and swarm size costs one search and buys a row
that holds up across the whole held-out grid, where a single-point row does not.
Neither that objective nor the doubled budget it came with was ablated, so which
of the two does the work is not known.

\emph{Report the baseline as a region.} The behavioural conclusion reproduces in
every seed draw and the returned parameters do not. A paper quoting a single
searched parameter vector without a seed study is quoting one sample from a
distribution it has not measured.

\emph{Match the estimator to the noise.} A search whose per-evaluation estimate
is noisier than the differences it must resolve fails in a way that looks exactly
like success: a plausible parameter vector, a good reported objective, no
held-out advantage. Two cheap diagnostics separate the cases: an honest re-score at high
replication, and a spread measurement at the hardest training condition.

\subsection{Plasticity, and what the axle reads}

The motivation for asking the capability question in hostile environments is the
observation, made as a position rather than a result by \citet{hunt2020}, that
field swarms meet conditions their designers did not enumerate and that minimal
designs have no margin to absorb them. The answer here is mixed in the way that
position anticipates: on one axis the minimum is unmoved and the cost is
performance, on the other the cost is capability.

One connection to sensing theory is available and must be stated at the strength
the data supports, which is an observation and not a result. The exact
decomposition of the terrain effect shows the heading-rate residual is
$v \cdot \partial m/\partial n + \omega \cdot (m(\text{centre}) - 1)$: the first
term is a difference in traction across the axle, so the axle does read a
spatial gradient, in the sense that \citet{bergpurcell1977} distinguish spatial
from temporal comparison at small scales. But the mean-traction term is roughly
ten times the gradient term at these constants, and no temporal-versus-spatial
comparison was run. It is an observation about what the mechanism decomposes
into, not evidence that a temporal comparator would do better.

\section{Limitations}
\label{sec:limitations}

The limitations below are the ones this study accepts rather than answers. Several
bound the claims above tightly enough that a reader could otherwise over-read
them.

\paragraph{Simulation only, and the reality gap is untested.}
Every number here comes from a purpose-built simulator. A second-tier replication
in an established simulator was specified in the project plan and never run, and
no hardware experiment exists. Contact is positional relaxation rather than
contact dynamics --- adequate for the comparisons made here, and precisely why a
second tier was planned. The automatic-design literature treats
the reality gap as a first-class design constraint
\citep{jakobi1995,francesca2014,ligot2018,birattari2019}; this study has taken
none of those steps. We claim simulation results and state plainly that transfer
is untested. Where a result is about protocol rather than about robots --- the winner's
curse, the noise floor, the seed scatter --- the reality gap is a weaker threat,
since those are properties of the search.

\paragraph{The dispersion normalisation is provisional.}
The dispersion metric follows the normalised second moment convention
\citep{grahamsloane1990}, under which a packed
cluster is designed to score near $1$; the clean baseline lands at $1.40$ and the
constant has never been checked against \citet{gauci2014}'s published figures.
Every comparison here is between cells sharing the normalisation, so no claim
depends on the absolute value --- but absolute dispersion should not be compared
across papers, swarm sizes, or body diameters.

\paragraph{Coverage.}
Almost every result is at $n = 20$; the exceptions are the clean-arena scaling
curve and the class-search grid at $n \in \{20, 50\}$. The main terrain results
use a single correlation length of $0.10$\,m. Peaks are located to a logarithmic
grid of ratio $1.39$, which for the turning-circle family is also the size of the
effect measured. The mission class is six points on two axes, so ``transfers''
means across start radius and swarm size only. Trial length was extended only in
the largest start-radius cells, and density is deliberately confounded with start
radius and swarm size wherever a sweep pins the start radius.

\paragraph{Rows that are not measurements of a capability.}
Every searched row is an upper bound from a diagonal-covariance optimiser at a
fixed budget: a stronger searcher could find more, which is the framework's own
rule rather than a defect. Every pursuer row above the baseline is hand-designed
and supports claims about spatial strategy, never about what a given number of
sensor states can do. The communication row was never run. Several geometry rows
are not buildable robots --- two put the wheel contacts outside the body shell ---
and exist only to break numerical confounds; the widest-axle row carries the two
largest $\ell/\lambda$ points, so that curve's shape beyond $\ell/\lambda = 1$
rests on it.

\paragraph{Confounds that remain inside the body-diameter result.}
Three lengths scale with the body and are not separated: the body as an
obstacle, the cluster link distance used by the metric, and the occlusion
footprint. Packing fraction also moves with the body at fixed start radius. The
log-log slope rests on two usable body sizes, and admitting the third --- which
reaches a cluster in at most $72\%$ of runs --- would halve it. The validity
window that excludes such rows was fixed after seeing the first sweep and is
labelled post hoc there; it was fixed in advance for the companion sweep and for
the body-diameter sweep, and the decision is read from those.

\paragraph{Saturation in the pursuer results.}
$8.8\%$ of runs in the first pursuer sweep and $29\%$ in the dispersive sweep
ended in a wipeout, concentrated exactly where the claims are strongest; the
survival collapses toward zero there --- 0.0199 per robot at the two longest
ranges --- and time-to-wipeout is what separates the rows. At the perfect-perception cell the base-task axis is not
estimable for three of the four aggregating rows. The pursuer sweeps use one
pursuit speed ratio, one targeting rule, one search strategy, one pursuer and
one trial length; two of three implemented targeting rules and one of two search
strategies were never swept.

\paragraph{Protocol deviations we must disclose.}
Four deviations from the study's own pre-registered statistical plan are on
record. (i) The planned Friedman test with post-hoc comparison across capability
rows was never run; the wrapper exists and raises rather than silently
substituting a weaker test, and every comparison in this paper is instead
pairwise with bootstrap or Wilson intervals, or paired by run index.
(ii) A single decision threshold $T$ was never pre-registered; the figures show
several contours ($T = 0.7, 0.8, 0.9$) and the text reads results off intervals.
The contours show how an answer depends on the bar; they are not a substitute
for having declared one. (iii) The plan asks for agreement between two
structurally different optimisers; the agreement check was run with three
independent seeds of the same optimiser, which cannot detect a bias common to
the diagonal-covariance family. (iv) The second experimental tier was never run.

\paragraph{Recorded and unexplained.}
Three observations are recorded without explanation and not claimed: the small
reach dip at $n = 3$ and $n = 4$, for which the simplest account is combinatorial
rather than deadlock; the scalar-field row's $\sim\!2\%$ improvement over flat
ground; and occlusion, run once and never revisited. The anisotropy result the
project set out to prove as a theorem was narrowed by measurement and never
attempted analytically.

\section{Conclusion}
\label{sec:conclusion}

We fixed a behaviour, parameterised the environment, and asked how the minimum
capability moves. Along a passive hostility axis it does not move. Two sensor
states, no memory and no arithmetic suffice over the whole swept terrain range;
small perturbations do not help, across three nulls; an added terrain state does
not help either, even from a warm start; and terrain acts as a performance tax of
which $96.9\%$ is objective mismatch and $3.1\%$ terrain at the tuned start
radius, $99.7\%$ against $0.3\%$ at twice it. The mechanism is per-wheel
curvature, not speed heterogeneity. How bad the worst case gets scales with the
turning circle; \emph{where} the worst case sits does not, and the claim that it
tracks the robot's body diameter instead is reported at SUGGESTED strength with
its hedge.

Along an active hostility axis the minimum does move, and the capability that
matters is the one that tells a pursuer from a neighbour --- among the
hand-designed rows tried\,\ddg, so a lower bound on what $S = 3$ offers rather
than a measurement of it. Confusion acts through aggregation, shown by a
dispersive control with identical sensing that gains
$1.20$--$1.13$ where aggregating rows gain $3.05$--$4.10$. Because that control is
nearly insensitive to confusion and the aggregating rows are not, \textbf{which
strategy wins at a given pursuer range is a function of $\kappa$, not of range
alone}, and no claim of the form ``aggregation wins inside this radius'' is well
formed without one.

The methodological findings are the ones we would most want carried forward. A
controller searched at one operating point can be worse than the published
constants at another, and a longer trial does not fix it. Searching over a mission
class produces a controller that transfers, though the ingredients that made it
work were not ablated. A search whose per-evaluation noise is nearly six times its
signal returns a plausible controller with a good reported objective and no real
advantage. And three optimiser seeds of that class search all beat the enumerated
reference while returning parameters that scatter by $16.9\%$ in turning circle:
the reproducible object is a region of behaviour, not a point in parameter
space.

\paragraph{Follow-up.}
The experimental record for this paper is frozen. The following are recorded as
follow-up work and were deliberately \emph{not} run for this paper: a class
search with more runs per condition, or a variance-stabilising statistic in
place of the per-condition median, to put the search above its own noise floor;
a sweep separating the three body-scaled lengths, most cheaply by varying the
metric's cluster link distance at a fixed body; a decoupling of the sensor
radius from the body; and anything with a finite-range sensor. Beyond those, the
record leaves open a searched dispersive row, a denser pursuer-range axis to
locate the crossover at two start radii, the communication row with an explicit
delivery model, a search budget scaled with dimension rather than warm started,
a second structurally different optimiser, the turning circle swept as a dial in
its own right to turn the gathering-rate-against-holding-quality explanation
into a measurement, whether the collapse onto axle-over-correlation-length
survives outside one decade, and the unexplained scalar-field improvement.

\section*{Disclosure}
\addcontentsline{toc}{section}{Disclosure}

\begin{framed}
\noindent\textbf{This document is an AI-generated draft.} It was written by
Claude (Anthropic) from a structured summary of the project's documented
findings, as a comparison baseline for a paper the repository owner is writing
by hand. There is no human author. It has not been reviewed, verified or
endorsed by any person, and it is not for submission, citation, or circulation.
Every number in it should be checked against the project's own records before
being relied upon. The disclosure markers throughout this document --- the title
prefix, the author block, the first sentence of the abstract, the running header
and footer, the page watermark, the marker on every figure caption, and this
paragraph --- are mandatory and must not be removed or weakened.
\end{framed}

\appendix

\section{Parameter tables}
\label{app:params}

\begin{table}[h]
\centering
\caption{Robot and simulator constants. The axle length is corrected from the
$5.3$\,cm value that circulates in secondary sources; $5.1$\,cm is what
reproduces the published derived quantities. The three derived quantities are
pinned by a regression test.}
\small
\begin{tabular}{lll}
\toprule
quantity & value & note \\
\midrule
body diameter & 7.4\,cm & disk, mass 152\,g \\
inter-wheel distance (axle) & 5.1\,cm & corrected from 5.3\,cm \\
wheel speed limit & $\pm 12.8$\,cm/s & per wheel \\
control cycle & 0.1\,s & physics substepped $10\times$ per cycle \\
sensor & zero-width, infinite-range ray & binary line of sight \\
state-0 turn radius $R_0$ & 14.45\,cm & reproduced to $5\times10^{-5}$\,m \\
$\omega_0$ & $-0.75$\,rad/s & reproduced to $5\times10^{-3}$\,rad/s \\
$\omega_1$ & $-5.02$\,rad/s & reproduced to $5\times10^{-3}$\,rad/s \\
start radius at $n = 20$ & 0.74\,m & uniform in a disc \\
contact resolution & positional relaxation & 32 passes, $0.01\%$ residual overlap \\
\bottomrule
\end{tabular}
\end{table}

\begin{table}[h]
\centering
\caption{Terrain model. The field is lattice value noise with quintic
interpolation; $\lambda$ is a lattice spacing rather than a Gaussian correlation
length, and the field's amplitude distribution is bell-shaped and concentrated
near zero rather than uniform.}
\small
\begin{tabular}{lll}
\toprule
quantity & value(s) & note \\
\midrule
traction amplitude $\theta_m$ & $0 \le \theta_m \le 0.9$ & cut below $1/\max|f| = 1.0$ \\
correlation length $\lambda$ & 0.10\,m (main); & lattice spacing \\
 & 0.02--0.386\,m (sweeps) & \\
field median $|f|$ & 0.3476 & $\max$ 0.9998, p99 0.9332 \\
traction floor & 0.05 & verified never to bind in range \\
slope term $g_{\mathrm{eff}}$ & 0.4 & uphill costs speed (ADR 0002) \\
per-wheel slip noise & 0 & implemented, unused in every sweep here \\
terrain bit threshold & $\theta_m \times 0.3476$ & held to a 0.4--0.6 split by test \\
\bottomrule
\end{tabular}
\end{table}

\begin{table}[h]
\centering
\caption{Pursuer model. Sensing range is quoted in metres and in start radii
throughout; $r_p \ge R$ is the perfect-perception corner rather than the hard end
of a difficulty axis.}
\small
\begin{tabular}{lll}
\toprule
dial & value(s) & definition \\
\midrule
speed ratio $\rho$ & 1.5 & pursuer speed / robot maximum \\
sensing range $r_p$ & 0.10, 0.20, 0.35, 0.60, 1.00\,m & omnidirectional, not a ray \\
$r_p / R$ & 0.14, 0.27, 0.47, 0.81, 1.35 & $R = 0.74$\,m at $n = 20$ \\
confusion $\kappa$ & 0, 0.5, 1.0, 2.5, 5.0 (3.0 in Pareto cells) & $p_{\mathrm{lock}} = 1/(1 + \kappa n_{\mathrm{local}})$ \\
confusion radius & 0.5\,m & window for $n_{\mathrm{local}}$ \\
handling time $h$ & 5.0\,s (first sweep) & --- \\
 & 0.39\,s and 1.93\,s (later sweeps) & $1\times$ and $5\times$ the 0.385\,s neighbour transit \\
targeting & nearest & two alternatives implemented, unswept \\
search & correlated random walk, 1.0\,rad/s & spiral implemented, unswept \\
capture distance & 0.08\,m & centre to centre \\
pursuer count & 1 & \\
trial length & 120\,s & \\
\bottomrule
\end{tabular}
\end{table}

\begin{table}[h]
\centering
\caption{Aggregation controller constants. Wheel speeds are fractions of the
$\pm 12.8$\,cm/s limit, given as (left, right) per sensor state. \dg\ searched,
not exhaustive --- an upper bound. \ddg\ hand-designed, not searched --- no
evidence about the capability itself.}
\small
\begin{tabular}{llllll}
\toprule
row & origin & state 0 & state 1 & $R_0$ & mark \\
\midrule
S2-gauci & exhaustive grid & $(-0.7000, -1.0000)$ & $(+1.0000, -1.0000)$ & 14.45\,cm & tight \\
S2-flat & searched, $\theta_m = 0$ & $(-0.3289, -0.8923)$ & $(+0.9983, -0.6589)$ & 5.53\,cm & \dg \\
S2-rough & searched, $\theta_m = 0.9$ & $(-0.2852, -0.9495)$ & $(+0.9354, -0.2262)$ & 4.74\,cm & \dg \\
S2-class-flat & class, $\theta_m = 0$ & $(-0.4330, -0.8820)$ & $(+0.7924, -0.8865)$ & 7.47\,cm & \dg \\
S2-class-rough & class, $\theta_m = 0.9$ & $(-0.2491, -0.9093)$ & $(+0.8620, -0.6283)$ & 4.47\,cm & \dg \\
S2-class-rough-$\tau$ & class, $\tau = 3600$\,s & $(-0.2367, -0.8740)$ & $(+0.7398, -0.2404)$ & 4.44\,cm & \dg \\
S4-composite & hand-built switch & Gauci on bit 0 & S2-rough on bit 1 & 14.45\,cm & \ddg \\
\bottomrule
\end{tabular}
\end{table}

\begin{table}[h]
\centering
\caption{The three optimiser seeds of the flat class search, identical in every
respect but the optimiser seed. The reported objective is what each search
believed it had achieved; the honest re-score evaluates the same constants at
100 runs per condition. All three are \dg.}
\small
\begin{tabular}{llllll}
\toprule
row & state 0 & state 1 & $R_0$ & reported & re-scored (gap) \\
\midrule
seed 1 & $(-0.4330, -0.8820)$ & $(+0.7924, -0.8865)$ & 7.47\,cm & 1.2143 & 1.2338 ($+0.0194$) \\
seed 2 & $(-0.5426, -0.9800)$ & $(+0.9467, -0.8113)$ & 8.88\,cm & 1.1927 & 1.2246 ($+0.0318$) \\
seed 3 & $(-0.4097, -0.7527)$ & $(+0.7326, -0.6060)$ & 8.64\,cm & 1.2120 & 1.2301 ($+0.0181$) \\
S2-gauci & $(-0.7000, -1.0000)$ & $(+1.0000, -1.0000)$ & 14.45\,cm & --- & 1.2905 \\
\bottomrule
\end{tabular}
\end{table}

\clearpage

\section{Revisions to the baseline during the study}
\label{app:corrections}

Fifteen corrections were made to the study's own design or to its reading of the
literature while it ran. They are listed because several of them changed results
that had already been recorded, and because a reader who knows only the final
design cannot tell which choices were forced.

\begin{enumerate}
\item \textbf{The two-robot proof of \citet{gauci2014} is unsound.} Theorem~3
claims two robots always aggregate; \citet{steinberg2025} identify the
unsound assumption and report a $4.24\%$ two-robot failure rate for the same
controller. A validation gate built on the theorem failed for two weeks before
the cause was found in the literature rather than in the code.

\item \textbf{An impossibility result covers the whole capability class.} For any
bimodal controller of this family there exists a swarm size and an initial state
for which aggregation does not occur, which makes the object of study
$c^*(\theta, n)$ rather than $c^*(\theta)$.

\item \textbf{The deadlock threshold is $n > 3$, and the quoted mechanism
matters.} \citet{daymude2021} establish deadlock for $n > 3$ under uniform
deterministic motion, with noise acting to break the precise balancing of
forces. This is the source of the noise-helps hypothesis and the reason it does
not transfer to a continuous model.

\item \textbf{Simulation parameters must be stated, and one was wrong.} The
inter-wheel distance is $5.1$\,cm, not the $5.3$\,cm carried in the initial
design; the derived quantities do not reproduce at $5.3$.

\item \textbf{The noise-helps hypothesis was withdrawn.} Three independent tests
returned nulls, including a 200-runs-per-cell neighbourhood of zero in which no
cell of 23 has an interval disjoint from baseline. The mechanism it was built on
is contact-deadlock breaking in a discrete model, and this continuous model
shows no such deadlock at $n = 20$.

\item \textbf{The aggregation threshold belongs on dispersion, not on a cluster
count.} Dispersion is invariant to the link distance used to define a cluster
($1.401$ at every value from 2.2 to 6.0 body radii) while the single-cluster
share moves from $0.43$ to $0.98$ over the same range.

\item \textbf{Handling time is required, not optional.} Without it a packed
cluster is captured wholesale in seconds and dilution cannot exist as an effect.

\item \textbf{A dispersive control row is required.} Without a row that has the
same sensing and the opposite spatial strategy, no confusion result can be
attributed to aggregation rather than to the sensor.

\item \textbf{The terrain model was replaced, and the replacement has a control.}
A scalar speed field leaves the instantaneous centre of rotation fixed, so it
cannot produce the effect under study; the per-wheel field can, and the scalar
field is retained as the control that demonstrates it.

\item \textbf{Report terrain amplitude only up to $1/\max|f|$.} Above it a wheel
can stall and the robot pivots, which is a different dynamical regime. The
$\theta_m = 1.25$--$1.5$ band in an early figure was stall, and was removed from
the swept range rather than explained.

\item \textbf{Re-tuning does not solve terrain.} It recovers the loss, but the
recovery decomposes into $96.9\%$ objective-tuning and $3.1\%$ terrain-tuning at
the tuned start radius, $99.7\%$ against $0.3\%$ at twice it, and the
terrain-specific share reverses sign on flat ground.

\item \textbf{The pursuer results must be reported on two axes.} Survival and
base-task performance differ by two orders of magnitude between strategies, so a
single scalar score hides the trade rather than summarising it.

\item \textbf{The worst correlation length is not a ratio law.} Two controllers
differing threefold in turning circle share a worst $\lambda$ in metres, and
inside one family doubling the turning circle moves the peak by $1.39\times$
with an interval excluding $2.00$. The original ratio reading stands as measured
for one controller and loses its generalisation.

\item \textbf{The $\ell/\lambda$ collapse no longer needs its single-$\lambda$
caveat.} Measured across a fourfold range of $\lambda$, the maximum spread in
slope at matched $\ell/\lambda$ is $0.0201$ against a range of $0.86$ along the
trend.

\item \textbf{Search the mission class, not one arena --- and report the class.}
A controller searched at one operating point can be worse than the published
constants at another; the objective, the conditions and the per-condition
statistic are part of the result and must be reported with it.
\end{enumerate}

\clearpage

\section{Claims register}
\label{app:claims}

Grades are assigned as follows. \textbf{SUPPORTED}: disjoint confidence
intervals, a replicated measurement, or a pre-registered decision rule that
fired. \textbf{SUGGESTED}: overlapping intervals, a single geometry, a post hoc
window, a single optimiser seed, or an association without a controlled cause;
these are stated only with the hedge given. \textbf{NOT SUPPORTED}: a reading
the data once seemed to support and no longer does, listed so that it cannot
return.

\subsection*{Supported}

\begin{longtable}{p{0.06\textwidth}p{0.86\textwidth}}
\toprule
\# & claim \\
\midrule
\endhead
S1 & Along the terrain dial, over $0 \le \theta_m \le 0.9$, the capability minimum does not move: two sensor states, no memory and no arithmetic suffice at every point, at both start radii tested. \\
S2 & No four-state controller better than the searched two-state one was found at equal budget, including when warm started at the two-state optimum. \\
S3 & A hand-built four-state composite switching between the two two-state rows is worse than either at every $\theta_m$, and reproduces the published constants exactly on flat ground. \\
S4 & Terrain is a performance tax whose size depends mainly on how well matched the controller is: hold ratios $2.149$, $1.200$ and $1.099$ at $\theta_m = 0.9$, paired per run. The matched-controller cost is 10--20\%. \\
S5 & Of the published-to-rough-trained gap at $\theta_m = 0.9$, $96.9\%$ is objective-tuning and $3.1\%$ terrain-tuning at $R = 0.74$\,m, against $99.7\%$ / $0.3\%$ at $1.5$\,m. The terrain term is not distinguishable from zero at either radius. \\
S6 & A scalar speed field produces no degradation at all (pooled $0.980$) where the per-wheel field peaks at $2.149$. \\
S7 & The heading-rate residual is the first-order per-wheel traction effect: slope $0.9943$, $R^2 = 0.9999$. \\
S8 & The quality of that linearisation collapses on $\ell/\lambda$: maximum slope spread $0.0201$ across a fourfold range of $\lambda$, against a range of $0.86$ along the trend. \\
S9 & Peak \emph{location} does not scale with $R_0$: peak ratio $1.39$ \CI{0.72}{1.93} for a doubling of $R_0$, excluding $2.00$. \\
S9b & Peak \emph{magnitude} does, and the axle does not set it: $5\%$ spread across a fourfold axle sweep at fixed $R_0$, $560\%$ across a fourfold $R_0$ sweep at fixed axle. S9 and S9b are separate results and must not be stated as one. \\
S10 & Confusion acts through aggregation: $\kappa: 0 \to 5$ multiplies aggregating rows' survival by $3.05$ and $4.10$ while a dispersive control with identical sensing gains $1.20$ and $1.13$; every endpoint interval is disjoint. \\
S11 & With survival pooled over $\kappa$ at $h = 1.93$\,s, the matched pair crosses between $0.27R$ and $0.47R$: aggregating wins at $0.14R$, the two overlap at $0.27R$, and the dispersive row wins from $0.47R$ outward. At a fixed $\kappa = 3$ the ordering at $0.47R$ reverses (S12), so the $\kappa$ is part of the claim. \\
S12 & At $r_p = 0.47R$ \emph{and} $\kappa = 3$ the dispersive row is Pareto-dominated: all four aggregating rows beat it on both survival and the base task, disjointly on survival. \\
S13 & Survival and base task differ by two orders of magnitude between strategies, so the pursuer results must be reported on two axes. \\
S14 & A controller searched at one operating point can be worse than the published constants at another, and a longer trial does not fix it. \\
S15 & The class-searched flat row matches or beats the published constants on dispersion in every cell at matched trial length (10 W, 2 tied, 0 L) and loses reach only at the largest start radius under terrain. \\
S16 & A training objective is the minimum of a noisy sample: reported $1.2595$ with no held-out advantage, and reported $1.6009$ against an honest re-score of $2.6553$. \\
S17 & At the worst training condition the two-run estimate the search saw spans $2.79$ log units against a $0.48$ log unit signal --- noise $5.8\times$ the signal. \\
S18 & Reach and hold are different measurements: reach is $1.00$ in every terrain cell except one runaway corner, and dispersion is invariant to the link distance while cluster metrics are not. \\
S19 & The published derived quantities reproduce: $R_0 = 14.45$\,cm to $5\times10^{-5}$\,m, $\omega_0 = -0.75$ and $\omega_1 = -5.02$\,rad/s to $5\times10^{-3}$\,rad/s. \\
S20 & Correlated and i.i.d. dropout must be compared at matched \emph{realised} rate: at a nominal $0.6$ the realised rate is $0.89$, and most of the apparent difference disappears when binned by realised rate. \\
S21 & The noise-helps hypothesis is null across three independent tests, including a dense 200-runs-per-cell neighbourhood of zero in which 0 of 23 cells has an interval disjoint from baseline. \\
S22 & A class search reliably produces a controller that beats the enumerated reference --- three seeds score $9/2/1$, $10/0/2$, $10/1/1$ --- but the returned controllers scatter: $R_0$ spans $16.9\%$ and 4 of 12 held-out cells contain a disjoint pair, so no single row is \emph{the} baseline. \\
\bottomrule
\end{longtable}

\subsection*{Suggested, with the wording used}

\begin{longtable}{p{0.06\textwidth}p{0.86\textwidth}}
\toprule
\# & claim, as it may be stated \\
\midrule
\endhead
G1 & Peak $\lambda$ tracks body diameter over a twofold range ($1.01$ and $0.97$ diameters, log-log slope $0.948$ \CI{0.474}{1.897}), on the two body sizes that aggregate at this start radius. Three lengths scale with the body and are not separated. Including the smallest body lowers the slope to $0.474$. No mechanism is proposed. \\
G2 & Proportionality in $R_0$ is rejected; a weak sub-proportional dependence is not excluded --- the point estimate is $1.39$ with an interval spanning $0.72$ to $1.93$, and peaks are located only to a log grid of ratio $1.39$. \\
G3 & The flat class row is a better point on the rough objective than the rough-trained row ($1.6456$ against $2.6553$ at 100 runs per condition on the rough search's own training seeds), which shows the search did not find its own optimum. It does not locate that optimum. \\
G4 & One sensor state that distinguishes pursuer from robot is worth more than anything else on the capability axis \emph{among the hand-designed rows tried}; these are \ddg, so this is a lower bound on what three states offer, and the five-state row's poor showing is a hand-design artefact. \\
G5 & Time to first cluster is monotone in $R_0$ across the five rows measured ($140$, $245$, $280$, $330$, $340$\,s); $R_0$ was never swept as a dial in its own right, so this is an association across five controllers rather than a controlled measurement. \\
G6 & Neither class-search ingredient (geometric-mean objective, doubled budget) was ablated; the row that transfers used both. \\
G7 & More runs per condition would raise the class search above its noise floor: the remedy follows from the measurement but was not run. \\
G8 & A scalar speed field aggregates $\sim\!2\%$ better than flat ground (pooled $0.980$, an interval excluding $1.0$). Reported and not chased. \\
G9 & The $n = 3$ / $n = 4$ reach dip is recorded, not claimed; the simpler explanation is combinatorial, and distinguishing it from deadlock needs the configurations checked directly. \\
G10 & The residual decomposes into a gradient term and a mean-traction term, so the axle does read a spatial difference --- but the mean-traction term is roughly ten times the gradient term at these constants and no temporal-versus-spatial comparison was run. An observation, not a result. \\
\bottomrule
\end{longtable}

\subsection*{Not supported --- readings ruled out}

\begin{longtable}{p{0.06\textwidth}p{0.86\textwidth}}
\toprule
\# & the reading, and why it is dead \\
\midrule
\endhead
N1 & ``Small amounts of terrain or noise help aggregation.'' Three independent nulls; 0 of 23 cells disjoint from baseline in a dense neighbourhood of zero. \\
N2 & ``The terrain bit cuts degradation from $2.21$ to $1.22$.'' True only without the searched two-state control; with it the two-state row reaches $1.10$ and the four-state row is worse, with disjoint intervals. \\
N3 & ``More sensor states hurt.'' The five-state table is a hand-written guess; it says no good five-state controller was found by hand. \\
N4 & ``Re-tuning solves the terrain problem.'' $96.9\%$ of the recovery is objective-tuning, and the terrain share reverses sign on flat ground. \\
N5 & ``The re-searched controller is the better controller.'' It is better only near its training condition; at four times the training start radius it fails outright, and at $n = 50$ it is worse than the published constants on flat ground at every radius. \\
N6 & ``The terrain effect peaks at $\lambda/R_0 \approx 0.7$ for any controller.'' Two controllers differing threefold in $R_0$ peak at the same $\lambda$ in metres. \\
N7 & ``There is a trade-off between flat-trained and rough-trained controllers that a terrain bit could arbitrate.'' The crossing exists at the tuned start radius and is gone at twice it. \\
N8 & ``The rough class row failed because its training objective was truncated.'' Removing the truncation moved $R_0$ from $4.47$ to $4.44$\,cm and made the row slightly worse; the cause is the noise floor. \\
N9 & ``Correlated dropout is much worse than i.i.d.\ dropout.'' At matched realised rate most of the difference disappears. \\
N10 & ``Two robots always aggregate.'' The proof is unsound and the same controller fails on $4.24\%$ of two-robot trials. \\
N11 & ``The $\theta_m = 1.25$--$1.5$ band shows terrain destroying aggregation.'' That band is wheel stall, not terrain. \\
N12 & ``A wider sensor cone fixes the small-$n$ failure.'' It is not what the original work did, and under the corrected criterion widening the cone makes $n = 2$ worse (reach $0.90 \to 0.17$). \\
\bottomrule
\end{longtable}

\clearpage

%% Bibliography entries are transcribed from the project's literature record,
%% which stores author, venue and year but not always the full title or the
%% volume and page numbers.  Missing fields are marked rather than invented.

\begin{thebibliography}{99}

\bibitem[Graham and Sloane(1990)]{grahamsloane1990}
Graham and Sloane (1990).
\newblock Penny-packing and two-dimensional codes.
\newblock \emph{Discrete \& Computational Geometry} 5.
\newblock \textit{Title supplied externally in revision; verify before submission.}

\bibitem[Holling(1959)]{holling1959}
Holling (1959).
\newblock Some characteristics of simple types of predation and parasitism.
\newblock \emph{The Canadian Entomologist} 91.
\newblock \textit{Title supplied externally in revision; verify before submission.}

\bibitem[Berg and Purcell(1977)]{bergpurcell1977}
Berg and Purcell (1977).
\newblock Physics of chemoreception.
\newblock \emph{Biophysical Journal}.

\bibitem[Birattari et al.(2019)]{birattari2019}
Birattari et al. (2019).
\newblock Automatic off-line design of robot swarms: a manifesto.
\newblock \emph{Frontiers in Robotics and AI}.
\newblock \textit{Title supplied externally in revision; verify before submission.}

\bibitem[Brown et al.(2016)]{brown2016}
Brown, Turgut and Goodrich (2016).
\newblock Discovery and exploration of novel swarm behaviors given limited robot
capabilities.
\newblock \emph{Distributed Autonomous Robotic Systems (DARS)}, 2016; Springer
volume 2018.

\bibitem[Chung et al.(2011)]{chung2011}
Chung, Hollinger and Isler (2011).
\newblock Search and pursuit-evasion in mobile robotics: a survey.
\newblock \emph{Autonomous Robots}.

\bibitem[Daymude et al.(2021)]{daymude2021}
Daymude, Harasha, Richa and Yiu (2021).
\newblock Deadlock and noise in self-organized aggregation without computation.
\newblock \emph{Symposium on Stabilizing, Safety, and Security of Distributed
Systems (SSS)}.

\bibitem[Francesca et al.(2014)]{francesca2014}
Francesca et al. (2014).
\newblock AutoMoDe: a novel approach to the automatic design of control software
for robot swarms.
\newblock \emph{Swarm Intelligence} 8(2).
\newblock \textit{Title supplied externally in revision; verify before submission.}

\bibitem[Gauci et al.(2014a)]{gauci2014}
Gauci, Chen, Li, Dodd and Gro\ss\ (2014a).
\newblock Self-organized aggregation without computation.
\newblock \emph{The International Journal of Robotics Research}.

\bibitem[Gauci et al.(2014b)]{gauci2014aamas}
Gauci, Chen, Li, Dodd and Gro\ss\ (2014b).
\newblock Clustering objects with robots that do not compute.
\newblock \emph{International Conference on Autonomous Agents and Multi-Agent
Systems (AAMAS)}.

\bibitem[Hamann(2018)]{hamann2018}
Hamann (2018).
\newblock \emph{Swarm Robotics: A Formal Approach}.

\bibitem[Hunt(2020)]{hunt2020}
Hunt (2020).
\newblock Phenotypic plasticity provides a bioinspiration framework for minimal
field swarm robotics.
\newblock \emph{Frontiers in Robotics and AI}.
\newblock \textit{Title supplied externally in revision; verify before submission.}

\bibitem[Jakobi et al.(1995)]{jakobi1995}
Jakobi, Husbands and Harvey (1995).
\newblock Noise and the reality gap: the use of simulation in evolutionary
robotics.
\newblock \emph{ECAL}.
\newblock \textit{Title supplied externally in revision; verify before submission.}

\bibitem[Johnson and Brown(2015)]{johnson2015}
Johnson and Brown (2015).
\newblock Evolving and controlling perimeter, rendezvous, and foraging behaviors
in a computation-free robot swarm.
\newblock \emph{EAI BICT}.
\newblock \textit{Title supplied externally in revision; verify before submission.}

\bibitem[Ligot and Birattari(2018)]{ligot2018}
Ligot and Birattari (2018).
\newblock On mimicking the effects of the reality gap with simulation-only
experiments.
\newblock \emph{ANTS}.
\newblock \textit{Title supplied externally in revision; verify before submission.}

\bibitem[Olson et al.(2013)]{olson2013}
Olson, Hintze, Dyer, Knoester and Adami (2013).
\newblock Predator confusion is sufficient to evolve swarming behaviour.
\newblock \emph{Journal of the Royal Society Interface} 10(85).
\newblock \textit{Title supplied externally in revision; verify before submission.}

\bibitem[\"Ozdemir et al.(2017)]{ozdemir2017}
\"Ozdemir, Gauci and Gro\ss\ (2017).
\newblock Shepherding with robots that do not compute.
\newblock \emph{European Conference on Artificial Life (ECAL)}.

\bibitem[\"Ozdemir et al.(2018)]{ozdemir2018}
\"Ozdemir, Gauci, Bonnet and Gro\ss\ (2018).
\newblock Finding consensus without computation.
\newblock \emph{IEEE Robotics and Automation Letters}.

\bibitem[Ros and Hansen(2008)]{ros2008}
Ros and Hansen (2008).
\newblock A simple modification in CMA-ES achieving linear time and space
complexity.

\bibitem[Schmickl et al.(2008)]{schmickl2008}
Schmickl, Kernbach et al. (2008).
\newblock Get in touch: cooperative decision making based on robot-to-robot
collisions (\emph{Autonomous Agents and Multi-Agent Systems} 18(1), 2009); with
Kernbach et al., Re-embodiment of honeybee aggregation behavior in an artificial
micro-robotic system (\emph{Adaptive Behavior}, 2009). BEECLUST.
\newblock \textit{Title supplied externally in revision; verify before submission.}

\bibitem[Steinberg and Solovey(2025)]{steinberg2025}
Steinberg and Solovey (2025).
\newblock Impossibility of self-organized aggregation without computation.
\newblock arXiv:2501.00390, submitted 31 December 2024 --- the verified source,
and the one cited here. An IEEE 2025 venue is recorded in the project's
literature list and is \textit{unverified}.

\end{thebibliography}

\end{document}
